\chapter{Phương pháp đề xuất}
\label{Chapter3}

\indent Trên cơ sở phân tích lý thuyết và yêu cầu thực tiễn, đề xuất kiến trúc hệ thống hỗ trợ học tập thông minh dựa trên LightRAG \cite{guo2025lightrag}, một mã nguồn mở được phát triển từ GraphRAG \cite{edge2024graphrag} và LangGraph \cite{langgraph2024}.
\section{Quy trình xử lý dữ liệu}

\begin{figure}[H]
	\centering
	\resizebox{\textwidth}{!}{%
\definecolor{cInputFill}{HTML}{505050}
\definecolor{cInputDraw}{HTML}{000000}
\definecolor{cProcessFill}{HTML}{FFFFFF}
\definecolor{cProcessDraw}{HTML}{000000}
\definecolor{cLLMFill}{HTML}{FFFFFF}
\definecolor{cLLMDraw}{HTML}{000000}
\definecolor{cSHACLFill}{HTML}{FFFFFF}
\definecolor{cSHACLDraw}{HTML}{000000}
\definecolor{cStorageFill}{HTML}{707070}
\definecolor{cStorageDraw}{HTML}{000000}
\definecolor{cSkipFill}{HTML}{F5F5F5}
\definecolor{cSkipDraw}{HTML}{9E9E9E}
\definecolor{cSkipText}{HTML}{757575}
\tikzset{
  >=Stealth,
  every node/.style={font=\sffamily\small},
  inputDoc/.style={rectangle,rounded corners=3pt,draw=cInputDraw,very thick,
    fill=cInputFill!90,text=white,minimum width=52mm,minimum height=12mm,
    align=center,text width=50mm},
  processBox/.style={rectangle,rounded corners=3pt,draw=cProcessDraw!80!black,thick,
    fill=cProcessFill,text=black,minimum width=52mm,minimum height=14mm,
    align=center,text width=50mm},
  llmBox/.style={rectangle,rounded corners=3pt,draw=cLLMDraw,thick,
    fill=cLLMFill,text=black,minimum width=52mm,minimum height=14mm,
    align=center,text width=50mm},
  shaclBox/.style={rectangle,rounded corners=3pt,draw=cSHACLDraw,thick,
    fill=cSHACLFill,text=black,minimum width=58mm,minimum height=14mm,
    align=center,text width=56mm},
  skipBox/.style={rectangle,rounded corners=3pt,draw=cSkipDraw,thick,dashed,
    fill=cSkipFill,text=cSkipText,minimum width=42mm,minimum height=12mm,
    align=center,text width=40mm},
  storageBox/.style={cylinder,shape border rotate=90,aspect=0.3,draw=cStorageDraw,thick,
    fill=cStorageFill,text=white,minimum width=32mm,minimum height=14mm,
    align=center,text width=28mm,font=\sffamily\tiny},
  phaseBox/.style={rectangle,rounded corners=6pt,draw=black!55,thick,dashed,
    inner xsep=8mm,inner ysep=8mm},
  phaseTitle/.style={font=\sffamily\bfseries\small,anchor=south west},
  solidEdge/.style={-{Stealth[length=2.5mm]},thick},
  dottedEdge/.style={-{Stealth[length=2.5mm]},thick,dash pattern=on 2pt off 2pt},
  lbl/.style={font=\sffamily\tiny\itshape,fill=white,inner sep=1pt},
}
\begin{tikzpicture}
% ── Pre-phase: Input → Dedup → Parse ──────────────────────────────────────────
\node[inputDoc]   (input)   at ( 0,    0  ) {Input Documents / Files / URL Links};
\node[processBox] (dedup)   at ( 0,   -2.2) {Document Deduplication\\(Kiểm tra tài liệu trùng lặp bằng MD5 Hash)};
\node[skipBox]    (skipdoc) at ( 9.0, -2.2) {Bỏ qua\\(Tài liệu đã tồn tại trong KV Storage)};
\node[processBox] (parse)   at ( 0,   -4.5) {MinerU Parse\\(Bóc tách Layout phức tạp, Bảng, PDF, Toán học)};

\draw[solidEdge] (input) -- (dedup);
\draw[solidEdge] (dedup) -- (skipdoc) node[lbl,midway,above]{Đã tồn tại};
\draw[solidEdge] (dedup) -- (parse)   node[lbl,midway,right]{Tài liệu mới};

% ── Phase 1: Dual Chunking ────────────────────────────────────────────────────
\node[processBox] (secchunk) at (-5.5, -7.0) {Section Chunking\\(Chia theo Heading/Depth)};
\node[processBox] (stchunk)  at ( 5.5, -7.0) {Standard Chunking\\(Chia theo Token Size)};

\node[storageBox] (seckv)   at (-7.2, -10.5) {Section KV Storage};
\node[storageBox] (secvdb)  at (-3.8, -10.5) {Section Vector DB};
\node[storageBox] (txkv)    at ( 3.8, -10.5) {Text Chunks\\KV Storage};
\node[storageBox] (txvdb)   at ( 7.2, -10.5) {Text Vector DB};

% split from parse with offset to avoid overlap at root
\draw[solidEdge] ([xshift=-3mm]parse.south) -- ++(0,-0.7) -| (secchunk.north);
\draw[solidEdge] ([xshift= 3mm]parse.south) -- ++(0,-0.7) -| (stchunk.north);
% chunking → storage (offset starting points)
\draw[dottedEdge] ([xshift=-3mm]secchunk.south) -- ++(0,-0.9) -| (seckv.north);
\draw[solidEdge]  ([xshift= 3mm]secchunk.south) -- ++(0,-0.9) -| (secvdb.north)
    node[lbl,pos=0.75,right]{use EB Model};
\draw[dottedEdge] ([xshift= 3mm]stchunk.south)  -- ++(0,-0.9) -| (txvdb.north);
\draw[solidEdge]  ([xshift=-3mm]stchunk.south)  -- ++(0,-0.9) -| (txkv.north)
    node[lbl,pos=0.75,left]{use EB Model};

\begin{scope}[on background layer]
\node[phaseBox,fit=(secchunk)(stchunk)(seckv)(secvdb)(txkv)(txvdb),
      inner xsep=7mm,inner ysep=7mm] (phase1) {};
\node[phaseTitle,yshift=5mm] at (phase1.north west) {1. Phân mảnh Kép (Dual Chunking)};
\end{scope}

% ── Phase 2: Extraction, Validation & Resolution ──────────────────────────────
\node[llmBox]     (extract)    at (-7.0, -14.5) {LLM Extraction\\(Bóc tách Entities \& Relations)};
\node[processBox] (orphan)     at ( 0,   -14.5) {Orphan Filter\\(Loại bỏ Entity cô lập)};
\node[llmBox]     (audit)      at ( 7.0, -14.5) {LLM Audit / Critic\\(Đánh giá confidence\_score)};

\node[shaclBox]   (shacl)      at ( 0,   -18.2) {SHACL Validation\\(Kiểm tra cấu trúc, rules)};
\node[shaclBox]   (decision)   at ( 0,   -22.0) {Decision Gate\\(Chấp nhận nếu shacl\_pass == True AND score $\geq$ 0.8)};
\node[skipBox]    (drop)       at ( 9.0, -22.0) {Bỏ qua\\(Rejected)};

\node[llmBox]     (resolution) at ( 0,   -26.5) {Entity Resolution (Coreference)\\(Nhận diện từ viết tắt, từ đồng nghĩa\\$\rightarrow$ Đồng nhất Entity \& Nối lại Relation)};

% Phase 2 entry from txkv (route through gap between phase1 and phase2)
\draw[solidEdge] (txkv.south) -- ++(0,-1.2) -| (extract.north);
\node[lbl] at (1.2,-12.4) {use LLM};

\draw[solidEdge] (extract) -- (orphan);
\draw[solidEdge] (orphan)  -- (audit);
\draw[solidEdge] (audit.south) |- (shacl.east);
\draw[solidEdge] (shacl)   -- (decision);
\draw[solidEdge] (decision)   -- (resolution) node[lbl,midway,left]{Pass Filter};
\draw[dottedEdge](decision)   -- (drop)        node[lbl,midway,above]{Reject};

\begin{scope}[on background layer]
\node[phaseBox,fit=(extract)(audit)(shacl)(decision)(resolution)(drop),
      inner xsep=7mm,inner ysep=7mm] (phase2) {};
\node[phaseTitle,xshift=12mm,yshift=5mm] at (phase2.north west)
      {2. Trích xuất, Xác thực \& Phân giải (Extraction \& Resolution)};
\end{scope}

% ── Phase 3: Dedup, Embedding & Graph Update ──────────────────────────────────
\node[processBox] (dedupent) at (-5.5, -31.0) {Deduped Entities\\(use 'Set' dedupe sau khi Resolution)};
\node[processBox] (deduprel) at ( 5.5, -31.0) {Deduped Relations\\(use 'Set' dedupe sau khi nối Relation)};

\node[llmBox]     (updent)   at (-10.5,-35.0) {Update Entity Description\\(use LLM để tổng hợp mô tả)};
\node[processBox] (embedent) at (-5.5, -35.0) {Entity Embedding};
\node[processBox] (embedrel) at ( 5.5, -35.0) {Relation Embedding};
\node[llmBox]     (updrel)   at ( 10.5,-35.0) {Update Relation Description\\(use LLM để tổng hợp mô tả)};

\node[storageBox] (kg)       at (-10.5,-39.5) {Knowledge Graph\\(Redis)};
\node[storageBox] (entvdb)   at (-5.5, -39.5) {Entity Vector DB};
\node[storageBox] (relvdb)   at ( 5.5, -39.5) {Relation Vector DB};

% Phase 3 split from resolution (offset ±3mm to show two distinct arrows)
\draw[solidEdge] ([xshift=-3mm]resolution.south) |- ([yshift=10mm]dedupent.north) -- (dedupent.north);
\draw[solidEdge] ([xshift= 3mm]resolution.south) |- ([yshift=10mm]deduprel.north) -- (deduprel.north);

\draw[solidEdge] (dedupent) -- (embedent) node[lbl,midway,right]{use EB Model};
\draw[solidEdge] (dedupent.west) -| node[lbl,pos=0.25,above]{use LLM} (updent.north);
\draw[solidEdge] (deduprel) -- (embedrel) node[lbl,midway,left]{use EB Model};
\draw[solidEdge] (deduprel.east) -| node[lbl,pos=0.25,above]{use LLM} (updrel.north);

\draw[dottedEdge] (embedent) -- (entvdb);
\draw[dottedEdge] (embedrel) -- (relvdb);
\draw[solidEdge]  (updent)   -- (kg);
% updrel → kg: route completely below all storage boxes to avoid crossing them
\draw[solidEdge]  (updrel.south) -- ++(0,-5.8) -| (kg.south);

\begin{scope}[on background layer]
\node[phaseBox,fit=(dedupent)(deduprel)(updent)(updrel)(entvdb)(relvdb)(kg),
      inner xsep=7mm,inner ysep=7mm] (phase3) {};
\node[phaseTitle,yshift=5mm] at (phase3.north west) {3. Chuẩn hóa Set \& Cập nhật Đồ thị};
\end{scope}

\end{tikzpicture}%
	}%
	\caption{Quy trình tổng quát từ văn bản thô đến đồ thị tri thức}
	\label{fig:flow_data_kg_main}
\end{figure}

\indent Đây là quy trình xử lý thông tin, đóng vai trò "nguyên liệu đầu vào" cho toàn bộ hệ thống. Hệ thống không chỉ lưu trữ dữ liệu thô mà còn thực hiện các phép biến đổi vector hóa phức tạp. Quản lý nguồn dữ liệu đa dạng bao gồm hồ sơ người học và các học liệu phi cấu trúc như bài giảng, tài liệu tham khảo.

\indent Đồ thị tri thức (Knowledge Graph) \cite{knowledge_graphs}: đây là thành phần quan trọng nhất, lưu trữ tri thức dưới dạng cấu trúc mạng lưới. Các nút đại diện cho các khái niệm và các cạnh biểu diễn mối quan hệ sư phạm. Điều này giúp hệ thống hiểu được "vị trí" và "tầm quan trọng" của một khái niệm trong tổng thể chương trình học.

\subsection{Tiền xử lý dữ liệu tri thức}

\subsubsection{Hệ thống MinerU}
Cơ sở lựa chọn công cụ và vai trò trong hệ thống như trong môi trường học thuật, nguồn dữ liệu đầu vào thường rất đa dạng và phức tạp, bao gồm các tệp PDF được scan từ văn bản giấy, tài liệu chứa nhiều biểu đồ, hình ảnh, và đặc biệt là các công thức toán học. Các công cụ trích xuất văn bản truyền thống thường xuyên gặp thất bại trong việc xử lý các tài liệu này, dẫn đến việc mất thông tin quan trọng hoặc làm xáo trộn cấu trúc bài giảng.

\indent Để giải quyết triệt để thách thức này, hệ thống áp dụng công cụ MinerU \cite{mineru_2024} (từ phiên bản 3.2.3 trở lên) làm giải pháp lõi cho quy trình tiền xử lý dữ liệu. Đây là một công cụ trích xuất dữ liệu đa phương thức chất lượng cao, tích hợp các mô hình tiên tiến như Vision-Language Model (VLM) cho việc phân tích bố cục (Layout Analysis), mô hình OCR để nhận diện văn bản trên ảnh, và mô hình Mathematical Formula Recognition (MFR) để chuyển đổi công thức toán học sang chuẩn LaTeX. Sự kết hợp này giúp MinerU không chỉ trích xuất chính xác văn bản mà còn bảo toàn hoàn hảo cấu trúc logic của tài liệu gốc, đặc biệt phù hợp với các giáo trình, bài báo khoa học và tài liệu lưu trữ của trường đại học.

\indent Để đạt hiệu quả tối ưu trong việc bóc tách các tài liệu học thuật phức tạp, MinerU phiên bản 3.x+ tích hợp một chuỗi các mô hình học sâu chuyên biệt cho từng loại tác vụ trong luồng xử lý hybrid. Cơ chế này giúp phân rã và xử lý chính xác từng phần của trang văn bản thông qua các mô hình được liệt kê chi tiết trong Bảng~\ref{tab:mineru_models}.

\begin{table}[H]
	\centering
	\renewcommand{\arraystretch}{1.4}
	\caption{Bảng tổng hợp các mô hình thành phần trong MinerU (Version 3.x+)}
	\label{tab:mineru_models}
	\begin{tabular}{|>{\raggedright\arraybackslash}p{4.2cm}|>{\raggedright\arraybackslash}p{4.0cm}|>{\raggedright\arraybackslash}p{6.8cm}|}
		\hline
		\textbf{Loại tác vụ (Task Type)} & \textbf{Tên mô hình (Model Name)} & \textbf{Vai trò và đặc điểm} \\ \hline
		Phân tích bố cục (Layout Detection) & PPDocLayoutV2 & Quét toàn bộ trang, nhận diện cấu trúc và phân loại 24 loại vùng dữ liệu khác nhau (văn bản, tiêu đề, ảnh, bảng biểu...). \\ \hline
		Nhận diện công thức (Formula Rec) & UniMERNet (Mặc định) hoặc PP-FormulaNet Plus M & Dịch các vùng ảnh chứa công thức (cả inline và block) sang mã chuẩn LaTeX. Bản Plus M bổ sung tối ưu cho công thức kèm ký tự tiếng Trung. \\ \hline
		Nhận diện chữ (OCR) & PytorchPaddleOCR (paddleocr\_torch) & Trích xuất ký tự văn bản thuần với độ phủ lên tới 109 ngôn ngữ. \\ \hline
		Phân loại bảng (Table Cls) & PaddleTableCls (dựa trên PP-LCNet) & Phân loại xem bảng đó thuộc dạng có viền (wired) hay không viền (wireless) để điều phối qua mô hình cấu trúc phù hợp. \\ \hline
		Cấu trúc bảng (Table Parsing) & SlanetPlus (cho bảng không viền) hoặc Unet (cho bảng có viền) & Bóc tách cấu trúc hàng/cột, xử lý các ô gộp phức tạp và chuyển đổi thành định dạng HTML/Markdown. \\ \hline
		Phát hiện hướng (Orientation) & PaddleOrientationCls (PP-LCNet) & Tự động phát hiện và xoay thẳng lại các trang hoặc bảng biểu bị đặt ngược hoặc nghiêng góc. \\ \hline
	\end{tabular}
\end{table}

\indent Bên cạnh các mô hình chuyên biệt nêu trên, trong các phiên bản cập nhật mới nhất (như dòng MinerU 2.5 Pro trở lên), hệ thống hỗ trợ tích hợp trực tiếp mô hình ngôn ngữ lớn đa phương thức Vision-Language Model (VLM) đóng vai trò là động cơ trích xuất duy nhất hoặc phối hợp (Hybrid). Khi hoạt động trong vai trò VLM chính của đường ống (pipeline), mô hình MinerU 2.5 Pro VLM tiếp nhận trực tiếp ảnh chụp trang tài liệu để thực hiện đồng thời các tác vụ nhận diện bố cục, nhận diện chữ (OCR), và dịch công thức toán học/bảng biểu chỉ qua một lượt suy luận duy nhất (single-turn inference). Giải pháp này giúp giảm bớt sự phụ thuộc vào chuỗi nhiều mô hình nhỏ kế tiếp nhau, hạn chế tối đa sai số tích lũy giữa các bước phân tích bố cục và nhận diện chữ, qua đó duy trì cấu trúc tài liệu một cách tự nhiên và chính xác hơn.

Luồng xử lý tự động (Hybrid-Auto-Engine) Chế độ \texttt{-b hybrid\allowbreak-auto\allowbreak-engine} trong MinerU thực thi một quy trình xử lý khép kín, kết hợp giữa mô hình đa phương thức (VLM) để bóc tách cấu trúc và các mô hình truyền thống (OCR, MFR) để nhận dạng văn bản. Dưới đây là sơ đồ chi tiết về luồng hoạt động này:

\begin{figure}[H]
	\centering
	\begin{tikzpicture}[
    scale=0.5, transform shape,
    >=Stealth,
    startstop/.style={rectangle, rounded corners, minimum width=2.6cm, minimum height=0.9cm, text centered, draw=black, fill=gray!25, font=\bfseries},
    process/.style={rectangle, rounded corners, minimum width=3.2cm, minimum height=1cm, text centered, text width=3.2cm, draw=black, fill=white},
    decision/.style={diamond, aspect=2.3, minimum width=3cm, minimum height=1cm, text centered, text width=2.4cm, draw=black, fill=gray!10},
    arrow/.style={thick,->,>=Stealth,rounded corners=4pt},
    line/.style={thick,rounded corners=4pt}
]

% ====================================================
% CỘT 1 (Trục chính bên trái, x=0)
% ====================================================
\node (start) [startstop] at (0, 0) {Bắt đầu};
\node (init) [process, text width=4.5cm] at (0, -2.4) {Khởi tạo cấu hình \& Tham số};

\node (doctype) [decision, aspect=2.3, text width=2.6cm] at (0, -5.6) {Loại tài liệu?};
\node (img) [process, text width=2.6cm] at (-5.8, -5.6) {Hình ảnh};
\node (office) [process, text width=3.4cm] at (6.8, -5.6) {Xử lý tài liệu Office \\ (Word, Excel, PPT)};

\node (pdf) [decision, aspect=2.1, text width=2.6cm] at (0, -9.2) {Tình trạng PDF?};
\node (ocr_true) [process, text width=3.2cm] at (-4.5, -12.4) {Bật chế độ OCR \\ (Ảnh quét)};
\node (ocr_false) [process, text width=3.2cm] at (4.5, -12.4) {Tắt chế độ OCR \\ (Văn bản số)};

\node (init_json) [process, text width=4.2cm] at (0, -15.8) {Khởi tạo dữ liệu trung gian \\ (Middle JSON)};
\node (loop) [decision, aspect=2.3, text width=2.6cm] at (0, -19.2) {Xử lý theo đợt \\ (Batch Window)};

% Nhánh kết thúc đợt (Trục chính)
\node (post) [process, text width=4.2cm] at (0, -23.2) {Hậu xử lý dữ liệu \& Hoàn thiện JSON};
\node (export_md) [process, text width=4cm] at (0, -25.8) {Xuất ra file Markdown};
\node (clean_html) [process, text width=4.6cm] at (0, -28.6) {Chuyển ký tự HTML trong Markdown thành Markdown thuần};
\node (end) [startstop] at (0, -31.4) {Kết thúc};

% ====================================================
% SNAKE FLOW (Vòng lặp vắt ngang chữ S, tận dụng chiều rộng)
% ====================================================
% Hàng 1 (Vắt sang Phải)
\node (img_conv) [process, text width=4.2cm] at (9.5, -19.2) {Chuyển đổi thành hình ảnh};
\node (layout) [process, text width=4cm] at (15.0, -19.2) {Dự đoán bố cục (Layout)};
\node (filter) [process, text width=4.2cm] at (21.0, -19.2) {Lọc công thức trong Bảng / Ảnh / Biểu đồ};

% Hàng 2 (Cuộn xuống và vắt sang Trái)
\node (effort) [decision, aspect=2.3, text width=2.6cm] at (21.0, -23.0) {Mức độ phân tích?};
\node (effort_high) [process, text width=3.8cm] at (15.0, -23.0) {VLM Trích xuất toàn diện (2 bước)};
\node (effort_med) [process, text width=3.8cm] at (15.0, -25.8) {Xác định hướng bảng \& VLM Trích xuất};
\node (check_ocr) [decision, aspect=2.3, text width=2.6cm] at (9.5, -25.8) {Đang bật OCR?};

% Hàng 3 (Cuộn xuống và vắt ngược sang Phải)
\node (run_ocr) [process, text width=4cm] at (15.0, -29.0) {Trích xuất vùng văn bản bổ sung (Sidecars)};
\node (run_mfr) [process, text width=4cm] at (15.0, -31.8) {Nhận diện công thức \& Trích xuất văn bản};
\node (split) [process, text width=4.2cm] at (21.0, -30.4) {Tách và phân loại tiêu đề};
\node (save) [process, text width=4cm] at (21.0, -33.4) {Lưu kết quả đợt};

% ====================================================
% MŨI TÊN KẾT NỐI
% ====================================================
% --- TRỤC CHÍNH (CỘT 1) ---
\draw [arrow] (start) -- (init);
\draw [arrow] (init) -- (doctype);

\draw [arrow] (doctype.west) -- node[above] {Hình ảnh} (img.east);
\draw [arrow] (doctype.east) -- node[above] {Office} (office.west);
\draw [arrow] (doctype.south) -- node[right] {PDF} (pdf.north);

\draw [arrow] (office.south) |- (post.east);
\draw [arrow] (img.south) -- ++(0, -0.6) -| (ocr_true.north);
\draw [arrow] (pdf.west) -| node[above, near start] {Ảnh quét} (ocr_true.north);
\draw [arrow] (pdf.east) -| node[above, near start] {Văn bản số} (ocr_false.north);

\draw [line] (ocr_true.south) |- (0, -14.2);
\draw [line] (ocr_false.south) |- (0, -14.2);
\draw [arrow] (0, -14.2) -- (init_json.north);

\draw [arrow] (init_json.south) -- (loop.north);
\draw [arrow] (loop.south) -- node[left] {Hết đợt} (post.north);
\draw [arrow] (post.south) -- (export_md.north);
\draw [arrow] (export_md.south) -- (clean_html.north);
\draw [arrow] (clean_html.south) -- (end.north);

% --- KẾT NỐI SNAKE FLOW ---
% 1. Rẽ sang phải
\draw [arrow] (loop.east) -- node[above] {Còn trang} (img_conv.west);
\draw [arrow] (img_conv.east) -- (layout.west);
\draw [arrow] (layout.east) -- (filter.west);

% 2. Cuộn xuống và Rẽ sang trái
\draw [arrow] (filter.south) -- (effort.north);
\draw [arrow] (effort.west) -- node[above] {High} (effort_high.east);
\draw [arrow] (effort.south) |- node[below, near start] {Medium} (effort_med.east);

\draw [arrow] (effort_high.west) -| (check_ocr.north);
\draw [arrow] (effort_med.west) -- (check_ocr.east);

% 3. Cuộn xuống và Rẽ sang phải (Tránh đè chéo, gộp 2 đường vào Split)
\draw [arrow] (check_ocr.south) |- node[above, near end] {Có} (run_ocr.west);
\draw [arrow] (check_ocr.south) |- node[below, near end] {Không} (run_mfr.west);

\draw [line] (run_ocr.east) -| (18.2, -30.4);
\draw [line] (run_mfr.east) -| (18.2, -30.4);
\draw [arrow] (18.2, -30.4) -- (split.west);

% 4. Lưu và Lặp về
\draw [arrow] (split.south) -- (save.north);
\draw [arrow, rounded corners=8pt] (save.south) -- ++(0,-1.0) -| (-8.2, -19.2) -- (loop.west);

\end{tikzpicture}
	\caption{Quy trình Hybrid-Auto-Engine của MinerU\cite{mineru_2024}}
	\label{fig:hybrid_auto_engine}
\end{figure}




\indent Quy trình xử lý được thực hiện tuần tự qua các giai đoạn. Đầu tiên, hệ thống tiến hành phân loại và phân tích định dạng của tệp đầu vào để xác định là văn bản gốc hay bản quét. Dựa trên kết quả này, cơ chế nhận dạng ký tự quang học sẽ tự động được kích hoạt nếu cần thiết. Tiếp theo, hệ thống đánh giá tài nguyên phần cứng hiện có, cụ thể là dung lượng bộ nhớ đồ họa VRAM, để tính toán tỷ lệ phân lô phù hợp, đảm bảo tối ưu hóa hiệu suất mà không gây tràn bộ nhớ. Sau quá trình khởi tạo, tài liệu được phân chia thành các đoạn nhỏ để xử lý dần. Tại mỗi phân đoạn, hệ thống sử dụng một mô hình đa phương thức VLM mạnh mẽ để bóc tách cấu trúc tổng thể, nhận diện các khối văn bản, bảng biểu, hình ảnh và công thức. Đối với các vùng chứa hình ảnh và bảng biểu, hệ thống sẽ tiến hành làm mờ nhằm không gây nhiễu cho các mô hình phía sau. Đồng thời, các công thức toán học nằm trong dòng văn bản sẽ được mô hình MFR chuyên biệt nhận diện và dịch đổi sang định dạng mã chuẩn LaTeX. Các khối chữ chưa xử lý được chuyển qua mô hình nhận diện ký tự quang học OCR. Nhằm đảm bảo độ chính xác, các kết quả có độ tin cậy thấp sẽ bị hệ thống tự động lọc bỏ. Cuối cùng, ở giai đoạn tổng hợp, tọa độ của các khối thông tin được chuẩn hóa về một hệ quy chiếu chung. Tất cả các dữ liệu trích xuất từ mô hình đa phương thức, mô hình nhận diện chữ và mô hình công thức được gộp chung vào một cấu trúc dữ liệu trung gian. Quy trình này lặp lại cho đến khi xử lý hết các phân đoạn của tài liệu, sau đó hệ thống sẽ giải phóng tài nguyên bộ nhớ và hoàn tất việc tạo ra một văn bản đầu ra đồng nhất và giữ nguyên cấu trúc logic gốc. Nhờ đường ống Hybrid này, MinerU khắc phục hoàn toàn những nhược điểm của việc đọc tệp bằng \texttt{textract} hay \texttt{pypdf} thông thường, cung cấp dòng dữ liệu đầu vào sắc nét và có độ tin cậy tuyệt đối cho mô hình ngôn ngữ phía sau.

\subsection{Luồng logic xử lý và cơ chế kiểm tra trùng lặp}

\indent Một trong những thành phần then chốt giúp hệ thống vận hành hiệu quả và tiết kiệm chi phí là cơ chế kiểm soát dữ liệu tại cổng vào. Sau khi văn bản được trích xuất từ các nguồn đa dạng (PDF, Office, Web) và đưa về định dạng Markdown đồng nhất, hệ thống thực hiện quy trình kiểm tra trùng lặp (Deduplication) dựa trên mã băm nội dung.

\indent Cụ thể, hệ thống sử dụng thuật toán băm MD5 để tạo ra một "dấu vân tay" kỹ thuật số duy nhất cho toàn bộ nội dung văn bản. Mã băm này sau đó được đối chiếu với cơ sở dữ liệu tri thức hiện tại, dẫn đến hai kịch bản xử lý chính. Trong trường hợp trùng khớp hoàn toàn, nếu mã MD5 của tài liệu mới đã tồn tại, hệ thống sẽ ngay lập tức dừng quy trình xử lý nhằm ngăn chặn việc nạp lại dữ liệu dư thừa, bảo vệ tính toàn vẹn của đồ thị tri thức và tiết kiệm tối đa tài nguyên tính toán. Ngược lại, đối với trường hợp nội dung mới hoặc có sự thay đổi, tài liệu được coi là nguồn tri thức mới và được chuyển tiếp sang giai đoạn phân đoạn. Ngay cả khi tài liệu chỉ có một thay đổi nhỏ, hệ thống vẫn sẽ thực hiện lại quy trình để đảm bảo tính cập nhật, trong khi sự trùng lặp về mặt thực thể sẽ được tự động hợp nhất tại bước xây dựng đồ thị sau này.

\begin{figure}[H]
	\centering
	\begin{tikzpicture}[node distance=1.5cm, auto, scale=0.85, transform shape]
		% Nodes
		\node (in) [startstop] {Văn bản Markdown đồng nhất};
		\node (md5) [process, below=of in] {Tính toán mã băm MD5};
		\node (check) [decision, below=1.5cm of md5] {MD5 đã tồn tại?};
		
		\node (stop) [startstop, right=2cm of check] {Bỏ qua (Trùng lặp)};
		\node (chunk) [process, below=1.5cm of check] {Chuyển sang phân đoạn \\ (Chunking)};
		
		% Arrows
		\draw [arrow] (in) -- (md5);
		\draw [arrow] (md5) -- (check);
		\draw [arrow] (check) -- node[anchor=south, font=\footnotesize] {CÓ} (stop);
		\draw [arrow] (check) -- node[anchor=west, font=\footnotesize] {KHÔNG} (chunk);
	\end{tikzpicture}
	\caption{Luồng logic kiểm tra trùng lặp bằng cơ chế băm nội dung}
	\label{fig:flow_deduplication}
\end{figure}

\subsubsection{Hệ thống thu thập tri thức từ trang web}

\begin{figure}[H]
	\centering
	\begin{tikzpicture}[node distance=1.8cm, auto, scale=0.8, transform shape]
		\node (input) [startstop] {URL Mục tiêu};
		\node (dom) [process, below=of input] {Nhận diện cấu trúc DOM};
		\node (check) [decision, below=of dom, text width=3cm] {Đã tồn tại trong \\ Registry hệ thống?};
		
		\node (llm) [process, left=2.8cm of check, yshift=-1.5cm, text width=3.5cm] {Phân tích LLM \\ (Logic tô pô DOM)};
		\node (gen) [process, below=1cm of llm] {Tạo bộ chọn \\ CSS/JS tối ưu};
		\node (save) [process, below=1cm of gen] {Cập nhật Registry \\ (Lưu trữ bộ chọn)};
		
		\node (exe) [process, right=2.8cm of check, yshift=-1.5cm, text width=4cm] {Thực thi Scraper Engine \\ (Playwright \& Stealth)};
		\node (clean) [process, below=1.5cm of exe] {Làm sạch Markdown};
		\node (out) [startstop, below=4cm of check, yshift=-2cm] {Đầu ra Markdown cuối cùng};
		
		\draw [arrow] (input) -- (dom);
		\draw [arrow] (dom) -- (check);
		\draw [arrow] (check) -| node [anchor=south, xshift=-1.4cm, font=\footnotesize] {KHÔNG (Mới)} (llm);
		\draw [arrow] (check) -| node [anchor=south, xshift=1.4cm, font=\footnotesize] {CÓ (Đã biết)} (exe);
		\draw [arrow] (llm) -- (gen);
		\draw [arrow] (gen) -- (save);
		\draw [arrow] (save) |- ($(out.north west)!0.3!(out.north)$);
		\draw [arrow] (exe) -- (clean);
		\draw [arrow] (clean) |- ($(out.north east)!0.3!(out.north)$);
	\end{tikzpicture}
	\caption{Quy trình thu thập tri thức từ trang web với Smart Scraper}
	\label{fig:flow_scrap_web}
\end{figure}

\indent Hệ thống cào web thông minh hoạt động theo quy trình tự động hóa từ việc nhận diện cấu trúc trang đến trích xuất dữ liệu. Logic cốt lõi dựa trên công cụ Smart Scraper Engine với cơ chế vận hành qua bốn bước chính: lấy dấu vân tay, đối chiếu kho lưu trữ, phân tích cấu trúc bằng mô hình ngôn ngữ và thực thi bộ phân giải.

\indent Quy trình bắt đầu bằng việc tạo ra một bản tóm tắt cấu trúc cây DOM của trang web, tập trung vào các thuộc tính định danh như thẻ HTML, tên lớp và ID. Hệ thống trích xuất một bộ khung xương kỹ thuật của trang web để làm cơ sở cho việc nhận diện. Sau đó, bản tóm tắt này được đối chiếu với kho lưu trữ các bộ phân giải đã biết. Mỗi mục trong kho lưu trữ chứa các thông tin về dấu hiệu nhận biết, bộ chọn nội dung chính, các thành phần cần loại bỏ và chiến lược chờ đợi phù hợp cho từng loại công nghệ web như React, Vue, hay các hệ quản trị nội dung truyền thống.

\indent Trong trường hợp gặp một trang web mới không có trong kho lưu trữ, hệ thống sẽ gửi một phần cấu trúc DOM tiêu biểu cho mô hình ngôn ngữ lớn. Mô hình này có nhiệm vụ phân tích logic phân bổ nội dung và tự động sinh ra mã lệnh điều hướng. Mã lệnh này bao gồm các bộ chọn CSS tối ưu để trích xuất văn bản có ý nghĩa và các quy tắc loại bỏ những thành phần gây nhiễu như trình đơn điều hướng, biểu ngữ quảng cáo hoặc các đoạn mã kịch bản. Bộ phân giải mới sinh ra sẽ được kiểm tra và lưu lại vào kho lưu trữ để sử dụng cho các lần truy cập sau, giúp tối ưu hóa hiệu suất và giảm thiểu chi phí gọi API.

\indent Đối với các tên miền mặc định như fit.hcmus.edu.vn và hcmus.edu.vn, hệ thống đã được cấu hình sẵn các bộ phân giải chuyên biệt cho nền tảng DotNetNuke. Do đặc thù của nền tảng này thường chứa nhiều nội dung lồng ghép trong các khung phức tạp, hệ thống sử dụng các bộ chọn mục tiêu như .dnn\_contentPane và \#dnn\_ContentPane để định vị chính xác vùng chứa thông tin học thuật. Quy trình trích xuất văn bản tại đây không chỉ lấy chữ thuần túy mà còn thực hiện chuyển đổi các liên kết quan trọng sang định dạng Markdown để bảo toàn thông tin tham chiếu.

\indent Đặc biệt, đối với các trang có tính tương tác cao như danh mục môn học hoặc đề cương chi tiết, hệ thống áp dụng chế độ tương tác vét cạn. Chế độ này sử dụng thư viện Playwright để tự động thực hiện các hành động như nhấn vào các nút mở rộng, cuộn trang để tải thêm dữ liệu và chờ đợi các tiến trình AJAX hoàn tất. Điều này đảm bảo rằng toàn bộ tri thức về chương trình đào tạo, mô tả môn học và thông tin giảng viên được thu thập đầy đủ nhất, ngay cả khi các thông tin này nằm sau các thành phần giao diện động. Dữ liệu sau khi trích xuất sẽ được làm sạch và chuẩn hóa trước khi đưa vào các bước phân đoạn và xây dựng đồ thị tri thức tiếp theo.

\subsection{Chiến lược phân đoạn và chuẩn hóa văn bản}

\indent Quá trình xử lý khởi đầu bằng việc làm sạch dữ liệu đầu vào, loại bỏ các ký tự lỗi và đảm bảo văn bản tuân thủ chuẩn mã hóa quốc tế Unicode (UTF-8). Bước tiền xử lý này không chỉ chuẩn hóa định dạng mà còn loại bỏ các thẻ HTML dư thừa, khoảng trắng không cần thiết và các ký tự điều khiển có thể gây nhiễu cho mô hình ngôn ngữ ở các bước xử lý tiếp theo.

\indent Về mặt kỹ thuật, hệ thống áp dụng song song hai chiến lược phân đoạn chính nhằm phục vụ hai mục tiêu khác biệt trong kiến trúc tổng thể: chiến lược phân đoạn theo kích thước cố định phục vụ trích xuất đồ thị tri thức và chiến lược phân đoạn phân cấp đa phương thức phục vụ truy xuất ngữ cảnh chính xác.

\indent Chiến lược thứ nhất, phân đoạn theo kích thước cố định (fixed-size chunking), được thiết kế chuyên biệt để phục vụ quy trình trích xuất thực thể và quan hệ nhằm xây dựng đồ thị tri thức trong LightRAG \cite{lightrag}. Phương pháp này chia toàn bộ văn bản đã chuẩn hóa thành các đoạn có độ dài xác định trước được thiết lập cố định là 1200 token theo nghiên cứu và khuyến nghị của các tác giả, đi kèm với vùng chồng lấp (overlap window) cố định là 100 token giữa hai đoạn liên tiếp.

\indent Cơ chế chồng lấp này đóng vai trò như một ``cầu nối ngữ cảnh'' (context bridge): khi mô hình ngôn ngữ xử lý đoạn thứ $i$, mô hình đã có sẵn phần đuôi của đoạn thứ $i-1$ làm nền tảng, từ đó tránh được hiện tượng cắt đứt giữa chừng một câu hoặc một luận điểm quan trọng. Cụ thể, với kích thước đoạn mặc định là 1200 token và vùng chồng lấp 100 token, hệ thống đảm bảo rằng các thực thể và mối quan hệ nằm gần ranh giới giữa hai đoạn liên tiếp vẫn được mô hình ngôn ngữ lớn nhận diện đầy đủ trong ít nhất một đoạn.

\indent Chiến lược phân đoạn kích thước cố định được lựa chọn cho tác vụ trích xuất đồ thị do quy trình trích xuất thực thể và quan hệ trong LightRAG yêu cầu mỗi đoạn văn phải đủ dài để chứa ngữ cảnh phong phú, giúp mô hình ngôn ngữ lớn nhận diện chính xác các thực thể và suy luận được mối liên hệ giữa các thực thể này. Các đoạn quá ngắn sẽ thiếu ngữ cảnh, dẫn đến việc bỏ sót các quan hệ ngầm định giữa các câu không liền kề.

\indent Chiến lược thứ hai, phân đoạn phân cấp (hierarchical chunking), được thiết kế nhằm bảo toàn tối đa cấu trúc và ngữ nghĩa của tài liệu gốc, phục vụ cho việc truy xuất ngữ cảnh chính xác khi hệ thống nhận truy vấn từ người dùng. Chiến lược này tham khảo từ MultiDocFusion \cite{shin-etal-2025-multidocfusion}, một hệ thống phân đoạn phân cấp được đề xuất cho các tài liệu dài và phức tạp, đã chứng minh hiệu quả vượt trội trong việc cải thiện độ chính xác truy xuất từ 8--15\% so với các phương pháp phân đoạn truyền thống.

\indent Khác với phân đoạn kích thước cố định vốn xử lý văn bản như một chuỗi ký tự thuần túy, chiến lược phân cấp đa phương thức nhận diện rằng tài liệu học thuật không chỉ chứa văn bản tuyến tính mà còn bao gồm nhiều thành phần có cấu trúc và ngữ nghĩa khác nhau. Do đó, quy trình phân đoạn không thể chỉ dựa trên số lượng token mà cần xem xét ba loại cấu trúc: (i) cấu trúc phân cấp (structural), tức mối quan hệ cha-con giữa các tiêu đề, mục, đoạn văn; (ii) cấu trúc bố cục (layout), tức vị trí vật lý, cột, vùng bao đóng của các thành phần trên trang; và (iii) cấu trúc ngữ nghĩa (semantic), tức mối quan hệ về nghĩa giữa các đoạn nội dung dù không kề nhau về vị trí vật lý.

\indent Quy trình phân đoạn phân cấp đa phương thức trong hệ thống được triển khai qua bốn giai đoạn tuần tự, tương ứng với bốn bước chính trong kiến trúc MultiDocFusion \cite{shin-etal-2025-multidocfusion}:

\indent Giai đoạn 1: Phân tích cấu trúc tài liệu bằng thị giác máy tính (Vision-based Document Parsing). Hệ thống VMiner (đã trình bày ở mục trước) sử dụng các mô hình học sâu chuyên biệt để phân tích cấu trúc hình học của từng trang tài liệu. Quá trình này tự động phát hiện và phân loại các vùng nội dung thành các loại thành phần riêng biệt: vùng tiêu đề (heading), khối văn bản (text block), bảng biểu (table), hình ảnh và biểu đồ (figure/chart), danh sách (list), chú thích (footnote) và công thức toán học. Mỗi vùng được xác định bởi một hộp bao đóng (bounding box) chứa tọa độ chính xác trên trang, làm cơ sở cho việc sắp xếp thứ tự đọc hợp lý và phân loại nội dung ở các bước tiếp theo.

\indent Giai đoạn 2: Trích xuất nội dung và phân loại thành phần đa phương thức. Sau khi nhận diện các vùng nội dung, hệ thống tiến hành trích xuất nội dung tương ứng từ từng vùng. Đối với các vùng văn bản thuần túy, nội dung được trích xuất bằng các công cụ nhận dạng ký tự quang học (OCR) đã tối ưu hóa cho tiếng Việt. Đối với bảng biểu, hệ thống tái tạo cấu trúc dữ liệu dưới dạng Markdown hoặc HTML, bảo toàn nguyên vẹn thông tin hàng, cột và tiêu đề. Đối với hình ảnh và biểu đồ, hệ thống sinh mô tả văn bản tự động thông qua mô hình ngôn ngữ thị giác (Vision Language Model, VLM), kết hợp với chú thích (caption) gốc và văn bản ngữ cảnh xung quanh hình ảnh trong tài liệu. Bước phân loại này đảm bảo mỗi thành phần được xử lý bởi chiến lược phù hợp nhất với đặc trưng dữ liệu của thành phần đó.

\indent Giai đoạn 3: Xây dựng cây cấu trúc phân cấp (Hierarchical Structure Parsing). Đây là giai đoạn cốt lõi phân biệt chiến lược phân cấp với các phương pháp phân đoạn truyền thống. Hệ thống tái cấu trúc toàn bộ tài liệu thành một cây phân cấp (hierarchical tree) phản ánh tổ chức logic của nội dung, trong đó gốc cây là tài liệu, các nút trung gian là các mục (section) và tiểu mục (sub-section), còn các nút lá là các đơn vị nội dung cụ thể như đoạn văn, bảng, hình ảnh. Cấu trúc phân cấp này được suy luận từ hệ thống tiêu đề (heading hierarchy) trong tài liệu, cụ thể các cấp tiêu đề H1, H2, H3 tương ứng với các mức độ lồng nhau khác nhau trong cây. Trong trường hợp tài liệu không có tiêu đề rõ ràng hoặc cấu trúc tiêu đề bị nhiễu, hệ thống bổ sung cơ chế phân tích dựa trên mô hình ngôn ngữ lớn (tương tự kỹ thuật DSHP-LLM trong MultiDocFusion \cite{shin-etal-2025-multidocfusion}) để suy luận ranh giới ngữ nghĩa giữa các phần. Mỗi nút trong cây được gán một đường dẫn phân cấp (breadcrumb path) dạng: \texttt{[Tài liệu] > [Chương X] > [Mục X.Y] > [Tiểu mục X.Y.Z]}, cung cấp ngữ cảnh vị trí chính xác cho từng đơn vị nội dung.

\indent Giai đoạn 4: Gom nhóm và sinh phân đoạn theo chiến lược DFS (DFS-based Grouping). Giai đoạn cuối cùng thực hiện việc duyệt cây phân cấp theo chiều sâu (Depth-First Search) để gom nhóm các nút lá thành các phân đoạn hoàn chỉnh, tuân thủ giới hạn kích thước token tối đa đã được cấu hình. Thuật toán duyệt DFS đảm bảo rằng các nút con thuộc cùng một nút cha sẽ được ưu tiên gom vào cùng một phân đoạn, từ đó bảo toàn tính liên kết ngữ nghĩa theo cấu trúc phân cấp gốc. Cụ thể, quy tắc gom nhóm được áp dụng dựa trên từng loại thành phần nội dung.


	
Đối với thành phần là một đoạn văn ngắn (dưới 100 token), gộp với các đoạn văn tiếp theo thuộc cùng mục (section) cho đến khi đạt ngưỡng kích thước tối thiểu.

Đối với thành phần là một đoạn văn dài (vượt quá giới hạn token), áp dụng kỹ thuật tách nhận biết ranh giới câu (sentence-aware splitting) để tránh cắt ngang ý tưởng.
	
Đối với thành phần là bảng biểu, giữ nguyên toàn bộ bảng như một phân đoạn độc lập kèm theo chú thích bảng và ngữ cảnh mục chứa bảng. Đối với bảng quá lớn vượt giới hạn token, áp dụng chiến lược tách theo nhóm hàng kèm giữ nguyên hàng tiêu đề.
	
Đối với thành phần là hình ảnh hoặc biểu đồ, tạo phân đoạn hình ảnh chứa mô tả sinh bởi VLM, chú thích (caption) gốc và đoạn văn bản ngữ cảnh trước/sau hình ảnh, thiết lập liên kết chéo giữa các thành phần này.
	
Đối với thành phần là tài liệu trình chiếu (slide), giữ nguyên toàn bộ nội dung slide gồm tiêu đề, khung văn bản, danh sách và ghi chú diễn giả làm một phân đoạn; tách khi nội dung vượt giới hạn.


\indent Mỗi phân đoạn được tạo ra sẽ được nhúng (inject) thêm thông tin đường dẫn phân cấp (breadcrumb context) vào đầu nội dung, theo định dạng:

\begin{center}
	\texttt{[Document: \{tên\_tài\_liệu\}] > [\{đường\_dẫn\_mục\}] \{nội\_dung\_gốc\}}
\end{center}

\noindent Kỹ thuật nhúng ngữ cảnh phân cấp này giúp mô hình ngôn ngữ lớn có đủ thông tin định vị ngay cả khi phân đoạn bị tách khỏi tài liệu gốc. Breadcrumb được nhúng vào đầu nội dung nhưng không được tính vào số lượng token của phần nội dung chính, nhằm không ảnh hưởng đến chất lượng biểu diễn vector nhúng (embedding) của nội dung thực.

\begin{table}[H]
	\centering
	\caption{So sánh hai chiến lược phân đoạn trong hệ thống}
	\label{tab:so_sanh_chunking}

	\renewcommand{\arraystretch}{1.4}
	\begin{tabular}{|p{3.5cm}|p{5.5cm}|p{5.5cm}|}
		\hline
		Tiêu chí & Chiến lược 1: Kích thước cố định & Chiến lược 2: Phân cấp đa phương thức \\ \hline
		Mục đích chính & Trích xuất thực thể và quan hệ cho đồ thị tri thức & Truy xuất ngữ cảnh chính xác cho RAG \\ \hline
		Đơn vị phân đoạn & Cửa sổ token cố định (1200 token) & Đơn vị ngữ nghĩa tự nhiên (đoạn văn, bảng, hình) \\ \hline
		Bảo toàn cấu trúc & Không, xử lý văn bản thuần tuyến tính & Có, bảo toàn cấu trúc phân cấp gốc \\ \hline
		Ngữ cảnh vị trí & Dựa vào vùng chồng lấp & Breadcrumb path nhúng vào mỗi chunk \\ \hline
		Ưu điểm & Đơn giản, phân bổ tải đều, phù hợp cho trích xuất hàng loạt & Bảo toàn ngữ nghĩa, truy xuất chính xác, hỗ trợ đa phương thức \\ \hline
		Hạn chế & Có thể cắt ngang ý tưởng & Phức tạp hơn, phụ thuộc vào chất lượng phân tích cấu trúc \\ \hline
	\end{tabular}
\end{table}

\indent Sự phối hợp giữa hai chiến lược phân đoạn tạo nên một kiến trúc lưu trữ kép (dual-indexing architecture) phục vụ cho hai luồng truy xuất khác nhau trong hệ thống. Chiến lược phân đoạn kích thước cố định cung cấp nguyên liệu đầu vào cho quy trình xây dựng đồ thị tri thức, nơi các thực thể và quan hệ được trích xuất và lưu trữ trong cơ sở dữ liệu đồ thị, phục vụ cho các chế độ truy vấn cục bộ (local) và toàn cục (global) của LightRAG. Trong khi đó, chiến lược phân đoạn phân cấp đa phương thức cung cấp các phân đoạn có cấu trúc ngữ nghĩa chặt chẽ, được lưu trữ trong cơ sở dữ liệu vector kèm theo siêu dữ liệu phong phú (metadata), phục vụ cho việc truy xuất ngữ cảnh chính xác khi hệ thống cần bổ sung thông tin chi tiết từ tài liệu gốc. Kiến trúc kép này đảm bảo rằng hệ thống vừa có khả năng lập luận đa bước trên đồ thị tri thức, vừa có khả năng truy xuất ngữ cảnh tài liệu gốc một cách trung thực và đầy đủ. \\ \indent Phương pháp thứ nhất - phân đoạn theo kích thước cố định (fixed-size chunking) - chia văn bản thành các đoạn có độ dài xác định trước (thường từ 512 đến 1200 token), kèm theo một vùng chồng lấp (overlap window) chiếm khoảng 10–20\% kích thước đoạn giữa hai đoạn liên tiếp. Cơ chế chồng lấp này đóng vai trò như một ``cầu nối ngữ cảnh'': khi mô hình xử lý đoạn thứ i, nó đã có sẵn phần đuôi của đoạn thứ i-1 làm nền tảng, từ đó tránh được hiện tượng cắt đứt giữa chừng một câu hoặc một luận điểm quan trọng.

\indent Phương pháp thứ hai - phân đoạn theo cấu trúc tự nhiên (structure-aware chunking) - khai thác các dấu hiệu cấu trúc sẵn có trong tài liệu như tiêu đề, tiêu đề phụ, danh sách có đánh số, hay các đoạn văn được phân tách bởi ký tự xuống dòng kép. Hệ thống sử dụng các biểu thức chính quy (regex) kết hợp với phân tích cú pháp nhẹ (lightweight parsing) để nhận diện ranh giới tự nhiên này. Ưu điểm của phương pháp này là mỗi đoạn văn giữ nguyên tính nguyên vẹn ngữ nghĩa - ví dụ, toàn bộ một mục trong tài liệu kỹ thuật sẽ không bị tách ra giữa chừng.

\subsubsection{Lựa chọn mô hình nhúng (Embedding Model)}

\indent Để tối ưu hóa phân hệ truy xuất ngữ cảnh (Retriever) trong kiến trúc RAG và cải thiện trực tiếp chỉ số độ gọi bối cảnh (Context Recall), hệ thống đề xuất sử dụng mô hình nhúng mã nguồn mở intfloat/multilingual-e5-large-instruct được phát triển bởi Microsoft Research \cite{multilingual_e5}. Mô hình này sở hữu kiến trúc biểu diễn hai chiều Encoder sâu sắc, được huấn luyện dựa trên tập ngữ liệu khổng lồ đa ngôn ngữ và đã chứng minh năng lực vượt trội khi đạt điểm số cao nhất trong bảng xếp hạng năng lực xử lý tiếng Việt tiêu chuẩn VN-MTEB (Vietnamese Massive Text Embedding Benchmark) công bố tại hội nghị khoa học EACL 2026 \cite{vn_mteb_eacl}. Qua nghiên cứu thực nghiệm, sự lựa chọn này đem lại các giá trị tối ưu vượt bậc cùng chiến lược xử lý các hạn chế vật lý bao gồm các khía cạnh cụ thể.

\indent Về cơ chế học có chỉ dẫn theo tác vụ (Task-aware Instruction Tuning), điểm cải tiến cốt lõi của phiên bản Instruct so với các kiến trúc nhúng phẳng thông thường là khả năng hiểu sâu sắc mục đích nghiệp vụ thông qua câu lệnh chỉ dẫn hệ thống \cite{optimal_rag_textbook}. Trước khi thực hiện vector hóa, hệ thống sẽ tự động cấu hình chèn tiền tố ngữ cảnh định hướng: chèn \texttt{"query:"} vào trước câu hỏi thảo luận học vụ của sinh viên và chèn \texttt{"passage:"} vào trước các phân đoạn tài liệu bối cảnh thô. Cơ chế này giúp mô hình chủ động căn chỉnh ma trận biểu diễn không gian, kéo các vector câu hỏi và vector bối cảnh chứa câu trả lời học vụ lại gần nhau hơn, giải quyết bài toán truy xuất bối cảnh chính xác \cite{multilingual_e5}.

\indent Về khả năng căn chỉnh đa ngữ (Cross-lingual Alignment), do đặc thù của kho tri thức học tập ngành Công nghệ thông tin đan xen dày đặc thuật ngữ tiếng Anh và tiếng Việt, mô hình E5 có khả năng ánh xạ các thực thể ngữ nghĩa đồng dạng của hai ngôn ngữ về cùng các vùng lân cận trong không gian vector 1024 chiều \cite{multilingual_e5}. Điều này giúp hệ thống phản hồi chính xác về bối cảnh ngay cả khi sinh viên đặt câu hỏi bằng tiếng Việt nhưng tài liệu đề cương được lưu trữ lại bằng tiếng Anh \cite{multilingual_e5}.

\indent Về hạn chế kỹ thuật của giới hạn vật lý và giải pháp tối ưu, hạn chế lớn nhất của mô hình nhúng này là cửa sổ ngữ cảnh đầu vào (Context Window) bị giới hạn nghiêm ngặt ở mức 512 tokens \cite{multilingual_e5}. Nếu phân đoạn văn bản vượt quá giới hạn này, mô hình sẽ kích hoạt cơ chế tự động cắt tỉa phần đuôi, gây thất thoát nghiêm trọng thông tin cấu trúc môn học. Để xử lý triệt để hạn chế này, hệ thống LightRAG đã thu hẹp cấu hình phân đoạn tịnh tiến (Chunking), hạ kích thước mỗi mảnh (Chunk Size) từ 1200 tokens mặc định xuống mức 400–500 tokens, phối hợp với cửa sổ dịch chuyển chồng lấp 50 tokens. Chiến lược này giúp bối cảnh văn bản lọt gọn vào giới hạn vật lý của mô hình biểu diễn mà vẫn bảo toàn ngữ nghĩa.

\indent Về mặt kiến trúc vận hành thực tế hệ thống, mô hình nhúng mở multilingual-e5-large-instruct sở hữu dung lượng tham số ở mức trung bình (560 triệu tham số) \cite{multilingual_e5}, mở ra hai kịch bản triển khai linh hoạt đáp ứng hoàn hảo từng môi trường hạ tầng:
\clearpage
\subsection{Cơ chế trích xuất thực thể và xây dựng quan hệ}

\indent Sau khi văn bản được chia thành các phân đoạn (chunk), LightRAG thực hiện quy trình trích xuất tri thức để chuyển đổi dữ liệu phi cấu trúc thành các thành phần đồ thị tri thức tường minh. Khác với các hệ thống RAG truyền thống thường bỏ qua mối liên kết giữa các đoạn văn, cơ chế này tập trung vào việc nhận diện các thực thể (Entities) và thiết lập các quan hệ (Relationships) xuyên suốt tài liệu.

\subsubsection{Chiến lược trích xuất đa nhiệm trong một lần gọi (One-step Extraction)}

\indent Để tối ưu hóa chi phí tính toán và duy trì tính nhất quán ngữ nghĩa, LightRAG sử dụng một lời nhắc có cấu trúc (structured prompt) yêu cầu mô hình ngôn ngữ lớn (LLM) thực hiện đồng thời hai nhiệm vụ trích xuất trong một lần gọi API duy nhất. 

\indent Nhiệm vụ 1: Nhận diện thực thể (Entity Extraction). Mô hình được yêu cầu xác định tất cả các thực thể quan trọng xuất hiện trong chunk. Đối với mỗi thực thể, LLM phải trích xuất các thuộc tính cốt lõi bao gồm tên định danh (Entity Name), loại thực thể (Entity Type) như PERSON, ORGANIZATION, CONCEPT, EVENT,... và một đoạn văn bản ngắn mô tả (Entity Description) giải thích vai trò, đặc điểm và ngữ cảnh của thực thể đó. Khác với NER truyền thống \cite{ner_survey}, phương pháp này khai thác tối đa khả năng hiểu ngữ nghĩa sâu của LLM để tạo ra các mô tả tri thức phong phú.

\indent Nhiệm vụ 2: Trích xuất quan hệ (Relation Extraction). Khác với các hệ thống đồ thị tri thức truyền thống chỉ sử dụng cấu trúc bộ ba (triple), hệ thống mở rộng biểu diễn thành cấu trúc bộ tứ (quadruple) với các thuộc tính. Sau khi xác định các thực thể, mô hình thực hiện quét các cặp thực thể để tìm kiếm các mối liên kết có ý nghĩa. Một quadruple trích xuất bao gồm hai thực thể tham gia (Source \& Target), các từ khóa định nghĩa bản chất liên kết (Relationship Keywords) và mô tả chi tiết (Relationship Description) lý do tại sao hai thực thể này có liên quan đến nhau dựa trên bằng chứng trong văn bản.

\subsubsection{Cơ chế Gleaning: Tự chỉnh sửa và Thu nhặt bổ sung}

\indent Một trong những thách thức lớn nhất của việc trích xuất tri thức bằng LLM là hiện tượng bỏ sót thông tin do giới hạn về khả năng chú ý (attention) khi xử lý các đoạn văn phức tạp. LightRAG giải quyết vấn đề này thông qua kỹ thuật \textit{Gleaning} (thu nhặt bổ sung).

\indent Sau vòng trích xuất đầu tiên, hệ thống không kết thúc ngay mà thực hiện gửi lại danh sách tri thức vừa tìm được cho LLM kèm theo một yêu cầu phản biện: ``Dựa trên đoạn văn này và danh sách thực thể/quan hệ đã có, liệu còn bất kỳ chi tiết nào bị bỏ sót không?''. Nếu mô hình phát hiện thêm tri thức mới, mô hình sẽ trả về kết quả bổ sung để hệ thống hợp nhất vào kết quả ban đầu. Quy trình này lặp lại cho đến khi mô hình xác nhận không còn thông tin sót hoặc đạt giới hạn số vòng lặp cấu hình.

\begin{figure}[H]
	\centering
	\begin{tikzpicture}[node distance=1.6cm, auto, scale=0.85, transform shape]
		% Nodes
		\node (input) [startstop] {Phân đoạn văn bản (Chunk)};
		\node (prompt) [process, below=1.1cm of input] {Tạo Prompt trích xuất \\ (NER + RE Schema)};
		\node (llm1) [process, below=1.1cm of prompt] {LLM Trích xuất lần 1};
		
		\node (loop_box) [groupbox, below=1.3cm of llm1, minimum width=11cm, minimum height=4.8cm] (loop) {};
		\node [above right=1mm of loop.north west, anchor=south west, font=\small\bfseries] {Quy trình Gleaning (Thu nhặt bổ sung)};
		
		\node (check) [decision, yshift=-1.2cm] at (loop.north) {Còn tri thức sót?};
		\node (llm_extra) [process, below=0.8cm of check] {Trích xuất bổ sung};
		
		\node (final) [startstop, below=1.2cm of loop] {Danh sách Thực thể \& Quan hệ cuối cùng};
		
		% Arrows
		\draw [arrow] (input) -- (prompt);
		\draw [arrow] (prompt) -- (llm1);
		\draw [arrow] (llm1) -- (check);
		\draw [arrow] (check) -- node[anchor=west, font=\footnotesize] {CÓ} (llm_extra);
		
		\draw [arrow] (llm_extra.east) -- ++(1cm,0) |- (check.east);
		
		\draw [arrow] (check.west) -- ++(-1.5cm,0) node[anchor=south, xshift=0.8cm, font=\footnotesize] {KHÔNG} |- (final.west);
		
	\end{tikzpicture}
	\caption{Luồng logic quy trình trích xuất tri thức và cơ chế Gleaning}
	\label{fig:flow_extraction_gleaning}
\end{figure}

\indent Kết quả cuối cùng là một tập hợp các thực thể và quan hệ có cấu trúc, sẵn sàng cho bước hợp nhất đồ thị (Graph Merging). Cơ chế này đảm bảo rằng mỗi đoạn văn bản đều được khai thác triệt để tri thức, tạo tiền đề cho việc xây dựng một đồ thị tri thức dày đặc và chính xác.

\subsection{Luồng xử lý Prompt linh hoạt cho đa dạng nguồn dữ liệu}

\indent Một thách thức lớn trong việc xây dựng đồ thị tri thức cho môi trường học thuật là sự khác biệt sâu sắc về cấu trúc và ngôn từ giữa các nguồn tài liệu. Đề cương chi tiết môn học thường mang tính quy phạm, ngắn gọn và sử dụng dày đặc các mã ký hiệu viết tắt như G1, G2 (mục tiêu môn học), CLO (chuẩn đầu ra), hay các thuật ngữ quản lý đào tạo như tích lũy, tiên quyết. Trong khi đó, dữ liệu từ hệ thống website của khoa lại mang tính chất thông tin đại chúng, chứa nhiều tin tức, sự kiện và mô tả tổng quát. Một lời nhắc trích xuất tổng quát duy nhất khó có thể bao quát tốt cả hai dạng tài liệu này: lời nhắc tối ưu cho đề cương sẽ bỏ sót ngữ cảnh tin tức, còn lời nhắc tối ưu cho website lại không nhận diện được các mã ký hiệu chuyên ngành. Để giải quyết vấn đề này, hệ thống triển khai mô-đun luồng xử lý lời nhắc linh hoạt (Dynamic Prompt Flow) đặt tại \texttt{lightrag/prompt\_flow}, với cơ chế tự động điều phối lời nhắc theo miền tri thức (domain) của từng phân đoạn.

\indent Về mặt thiết kế, mô-đun áp dụng đồng thời ba mẫu thiết kế phần mềm: \textit{Registry}, \textit{Strategy} và \textit{Protocol}. Mỗi miền tri thức được hiện thực thành một mô-đun chiến lược (strategy) tuân thủ giao diện \texttt{EntityExtractionModule}, bắt buộc cung cấp đúng bốn hằng số lời nhắc: lời nhắc hệ thống (\texttt{ENTITY\_\allowbreak EXTRACTION\_\allowbreak SYSTEM\_\allowbreak PROMPT}), tập ví dụ mẫu (\texttt{ENTITY\_\allowbreak EXTRACTION\_\allowbreak EXAMPLES}), lời nhắc người dùng (\texttt{ENTITY\_\allowbreak EXTRACTION\_\allowbreak USER\_\allowbreak PROMPT}) và lời nhắc thu nhặt bổ sung (\texttt{ENTITY\_\allowbreak CONTINUE\_\allowbreak EXTRACTION\_\allowbreak USER\_\allowbreak PROMPT}). Bộ điều phối (\texttt{PromptDispatcher}) duy trì một sổ đăng ký (Registry) ánh xạ tên miền sang mô-đun tương ứng; khi gặp một nhãn miền chưa được đăng ký, hệ thống tự động lùi về miền mặc định \texttt{curriculum}. Thiết kế này cho phép bổ sung một miền tri thức mới mà không phải sửa đổi mã nguồn của tầng trích xuất, đảm bảo khả năng mở rộng.

\indent Hệ thống hiện hỗ trợ bốn miền tri thức đặc thù cho dữ liệu đào tạo bậc đại học. Miền \texttt{curriculum} dành cho chương trình đào tạo, tập trung vào ngành, khoa, khối kiến thức, môn học, mục tiêu chương trình (PO), chuẩn đầu ra (ELO), số tín chỉ và điều kiện tiên quyết. Miền \texttt{course} dành cho đề cương chi tiết học phần, ưu tiên các chuẩn đầu ra môn học (CLO), nội dung chương, hình thức kiểm tra và tài liệu tham khảo. Miền \texttt{teacher} dành cho danh sách nhân sự giảng dạy, phân biệt giảng viên cơ hữu, thỉnh giảng và cán bộ hỗ trợ cùng học hàm, học vị. Cuối cùng, miền \texttt{resolution} dành cho các văn bản pháp lý, hành chính thuần túy như nghị quyết hội đồng trường hay quy chế tổ chức.

\indent Việc xác định miền của một phân đoạn được thực hiện theo cơ chế hai tầng nhằm cân bằng giữa độ chính xác và khả năng chịu lỗi. Tầng thứ nhất sử dụng một mô hình ngôn ngữ lớn đóng vai trò bộ phân loại nhẹ (lightweight classifier): hệ thống chỉ gửi khoảng 1500 ký tự đầu của phân đoạn kèm theo một lời nhắc phân loại ngắn gọn, yêu cầu mô hình trả về đúng một nhãn miền. Do các văn bản giáo dục Việt Nam thường mở đầu bằng phần thủ tục hành chính (``QUYẾT ĐỊNH V/v Ban hành...'', ``Căn cứ...''), lời nhắc phân loại được bổ sung một quy tắc phủ quyết (critical override): nếu phần mở đầu hành chính này thực chất nhằm ban hành chương trình đào tạo hoặc đề cương môn học, văn bản phải được phân vào \texttt{curriculum} hoặc \texttt{course} thay vì bị nhầm thành \texttt{resolution}. Tầng thứ hai là một bộ phân loại dựa trên biểu thức chính quy và tính điểm theo từ khóa đặc trưng (\texttt{classify\_heuristic}), đóng vai trò phương án dự phòng khi mô hình ngôn ngữ không khả dụng hoặc trả về kết quả không hợp lệ.

\indent Sau khi xác định được miền, bộ điều phối nạp các kiểu thực thể và kiểu quan hệ đặc thù từ lớp bản thể học (ontology) của miền đó vào các vị trí giữ chỗ \texttt{\{entity\_types\}} và \texttt{\{relation\_types\}} trong lời nhắc. Ví dụ, miền \texttt{curriculum} định nghĩa các kiểu thực thể như \texttt{Major}, \texttt{KnowledgeBlock}, \texttt{Course}, \texttt{LearningOutcome} (ELO), \texttt{ProgramObjective} (PO) cùng các kiểu quan hệ như \texttt{PREREQUISITE\_OF} (học trước), \texttt{CONTAINS} (khối kiến thức chứa môn học) hay \texttt{MAPS\_TO} (ELO liên kết với PO); khi miền không khai báo bản thể học, hệ thống dùng quan hệ tổng quát \texttt{RELATES\_TO} làm mặc định. Việc áp đặt sẵn bộ kiểu chuyên biệt giúp định hướng mô hình ngôn ngữ tập trung vào tri thức có giá trị nhất của miền và duy trì tính nhất quán về định danh (Entity Canonicalization) ngay từ bước sơ khởi, từ đó giảm thiểu số lượng thực thể phân mảnh và tạo điều kiện cho việc hợp nhất các nút đồng nghĩa ở các bước sau. Cần lưu ý rằng ở giai đoạn hiện tại, việc đối chiếu kết quả trích xuất với bản thể học chỉ dừng ở mức ghi cảnh báo (warn-only), còn cơ chế loại bỏ tự động dữ liệu vi phạm được để ngỏ như một điểm mở rộng trong tương lai.

\indent Để kiểm soát chi phí token, mô-đun không chèn toàn bộ tập ví dụ mẫu vào mỗi lần gọi mà áp dụng cơ chế chọn ví dụ động (dynamic few-shot selection). Một bộ chỉ mục TF-IDF kết hợp độ tương đồng cosin được xây dựng thuần Python (không phụ thuộc thư viện ngoài) để xếp hạng các ví dụ theo mức độ liên quan với nội dung phân đoạn hiện tại, sau đó chỉ giữ lại hai ví dụ (top-k = 2) phù hợp nhất. Cách làm này giảm đáng kể lượng token tiêu thụ so với việc nhúng toàn bộ ví dụ, đồng thời vẫn bảo toàn chất lượng trích xuất nhờ luôn cung cấp đúng ngữ cảnh mẫu sát với phân đoạn. Trong trường hợp việc xếp hạng gặp lỗi, hệ thống tự động lùi về phương án trả về các ví dụ đầu danh sách, bảo đảm quy trình không bị gián đoạn. Hình~\ref{fig:flow_prompt_dispatch} mô tả tổng quan luồng điều phối lời nhắc từ phân đoạn đầu vào đến bộ lời nhắc trích xuất hoàn chỉnh.

\begin{figure}[H]
	\centering
	\begin{tikzpicture}[node distance=1.3cm, auto, scale=0.82, transform shape]
		% Nodes
		\node (in) [startstop, text width=5cm] {Đoạn đầu phân đoạn \\ \footnotesize($\le$ 1500 ký tự)};
		\node (dec) [decision, below=1.1cm of in, text width=2.8cm] {LLM khả dụng?};
		\node (llm) [process, below left=1.1cm and 0.3cm of dec, text width=3.6cm] {Phân loại bằng LLM \\ \footnotesize(quy tắc phủ quyết)};
		\node (heur) [process, below right=1.1cm and 0.3cm of dec, text width=3.6cm] {Heuristic regex \\ \footnotesize(\texttt{classify\_heuristic})};
		\node (label) [process, below=3.4cm of dec, text width=5cm] {Nhãn miền \\ \footnotesize(curriculum / course / teacher / resolution)};
		\node (resolve) [process, below=1cm of label, text width=6cm] {Tra Registry $\rightarrow$ Module miền \\ \footnotesize(mặc định: curriculum)};

		\node (build) [groupbox, below=1.2cm of resolve, minimum width=9.5cm, minimum height=2.6cm] (bbox) {};
		\node [above right=1mm of bbox.north west, anchor=south west, font=\small\bfseries] {Dựng lời nhắc trích xuất};
		\node (onto) [smallbox, text width=4cm] at ($(bbox.center)+(-2.4cm,0)$) {Chèn \texttt{entity\_types} / \texttt{relation\_types} từ Ontology};
		\node (fewshot) [smallbox, text width=4cm] at ($(bbox.center)+(2.4cm,0)$) {Chọn top-k=2 ví dụ \\ \footnotesize(TF-IDF + cosin)};

		\node (out) [startstop, below=1.2cm of bbox, text width=6cm] {Bộ lời nhắc (system + user + continue)};

		% Arrows
		\draw [arrow] (in) -- (dec);
		\draw [arrow] (dec) -| node[anchor=south, pos=0.25, font=\footnotesize] {Có} (llm);
		\draw [arrow] (dec) -| node[anchor=south, pos=0.25, font=\footnotesize] {Không} (heur);
		\draw [arrow] (llm) |- (label);
		\draw [arrow] (heur) |- (label);
		\draw [arrow] (label) -- (resolve);
		\draw [arrow] (resolve) -- (bbox);
		\draw [arrow] (bbox) -- (out);
	\end{tikzpicture}
	\caption{Luồng điều phối lời nhắc linh hoạt theo miền tri thức (Dynamic Prompt Flow)}
	\label{fig:flow_prompt_dispatch}
\end{figure}

\subsection{Kiểm chứng tri thức bằng chuẩn SHACL và đối chiếu bằng chứng xác định}

\indent Để đảm bảo tính toàn vẹn và chất lượng của dữ liệu tri thức trước khi đưa vào đồ thị chính thức, hệ thống tích hợp một bộ kiểm chứng lấy cảm hứng từ chuẩn SHACL (Shapes Constraint Language) \cite{shacl_w3c}. SHACL cung cấp một ngôn ngữ để mô tả các ràng buộc về cấu trúc và ngữ nghĩa của dữ liệu đồ thị, cho phép tự động hóa việc phát hiện các lỗi phổ biến trong kết quả trích xuất bằng mô hình ngôn ngữ lớn \cite{shacl_wikidata}. Toàn bộ logic kiểm chứng được hiện thực trong mô-đun \texttt{lightrag/shacl}, tổ chức thành một quy trình đánh giá nghiêm ngặt và đối chiếu bằng chứng (deterministic text-grounding) thực thi trên kết quả trích xuất của từng phân đoạn.

\indent Lớp bảo vệ thứ nhất (Guard 1) -- Khử trùng lặp. Hệ thống đối chiếu mã định danh phân đoạn (\texttt{chunk\_id}) với kho lưu trữ hiện có. Nếu phân đoạn đã tồn tại, hệ thống chuyển sang trạng thái nạp lại (\texttt{RE\_INGEST}) và điều chỉnh cách tính bậc kết nối của thực thể để tránh cộng dồn trùng lặp; ngược lại, phân đoạn mới được cho phép tiếp tục (\texttt{PROCEED}). Lớp này ngăn việc xử lý lại dữ liệu dư thừa và bảo vệ tính toàn vẹn của đồ thị.

\indent Lớp bảo vệ thứ hai (Guard 2) -- Kiểm tra cấu trúc và đối chiếu bằng chứng. Đây là trọng tâm của quy trình. Mỗi thực thể và quan hệ được kiểm tra qua một tập các định nghĩa hình dạng (Shapes). Đối với thực thể, các quy tắc bao gồm: bắt buộc có tên hợp lệ (\texttt{missing\_name}, mức lỗi), bắt buộc có kiểu xác định khác ``UNKNOWN'' (\texttt{missing\_type}, mức lỗi), bắt buộc có siêu dữ liệu nguồn gốc \texttt{source\_id} (\texttt{no\_provenance}, mức lỗi), và cảnh báo thực thể bị cô lập khi số kết nối nhỏ hơn ngưỡng tối thiểu (\texttt{isolated\_entity}). Đáng chú ý, số kết nối của một thực thể được tính bằng cách kết hợp bậc cục bộ (trong phạm vi phân đoạn) với bậc toàn cục truy vấn đồng thời từ cơ sở dữ liệu đồ thị, nhằm tránh báo cô lập sai đối với những thực thể đã có liên kết ở các phân đoạn trước. Đối với quan hệ, hệ thống kiểm tra sự tồn tại của thực thể nguồn và đích, phát hiện quan hệ tự trỏ (\texttt{self\_loop}), độ dài tối thiểu của mô tả (mặc định 15 ký tự) và dải giá trị hợp lệ của trọng số trong khoảng $[0, 1]$. Ngoài ra, một bộ chặn ngôn ngữ (language guard) sẽ từ chối mọi giá trị chứa ký tự nằm ngoài bảng chữ Latinh và tiếng Việt (\texttt{contains\_foreign\_script}), loại bỏ các trường hợp mô hình ngôn ngữ vô tình sinh ra ký tự ngoại lai. Mỗi vi phạm được gán một trong ba mức nghiêm trọng: lỗi (error), cảnh báo (warning) hoặc thông tin (info).

\indent Điểm khác biệt cốt lõi của bộ kiểm chứng là cơ chế đối chiếu bằng chứng văn bản một cách xác định (deterministic text-grounding). Thay vì hỏi mô hình ngôn ngữ ``thực thể này có trong văn bản không?'' rồi tin vào câu trả lời, hệ thống tự tính toán bằng chứng: trước tiên tìm khớp chính xác chuỗi con (điểm số 1.0), nếu không có thì dùng so khớp mờ (fuzzy matching qua \texttt{rapidfuzz} hoặc \texttt{difflib}) với ngưỡng chấp nhận 0.9. Kết quả trả về không phải một ý kiến mà là chính đoạn văn bản khớp được kèm theo vị trí ký tự trong phân đoạn -- đóng vai trò bằng chứng có thể kiểm tra lại. Đối với quan hệ, hệ thống yêu cầu cả hai đầu mút đều có bằng chứng và đồng xuất hiện trong cùng một cửa sổ 200 ký tự, qua đó củng cố độ tin cậy rằng quan hệ là có thật chứ không phải hai lần nhắc rời rạc.

\indent Quyết định cuối cùng cho mỗi thực thể và quan hệ được tổng hợp tại tầng ra quyết định thuần túy theo nguyên tắc phân cấp niềm tin: các kiểm tra xác định (cấu trúc và bằng chứng văn bản) giữ quyền quyết định tuyệt đối, thay vì hoàn toàn phụ thuộc vào điểm số tự tin của LLM phản biện. Cụ thể, nếu vi phạm cấu trúc ở mức lỗi, dữ liệu bị từ chối (\texttt{REJECTED}); nếu hợp lệ về cấu trúc và có bằng chứng văn bản xác định, dữ liệu được phê duyệt (\texttt{APPROVED}); nếu hợp lệ về cấu trúc nhưng thiếu bằng chứng (ví dụ suy luận ngầm ẩn), dữ liệu chỉ được đưa vào trạng thái chờ xem xét (\texttt{PENDING\_REVIEW}) khi độ tin cậy của LLM phản biện đạt ngưỡng 0.8. Nghĩa là, LLM chỉ được phép đẩy các trường hợp ranh giới vào diện xem xét, chứ không bao giờ được quyền tự ý chấp nhận một tri thức vô căn cứ.

\indent Kết quả kiểm chứng điều hướng trực tiếp quá trình hợp nhất dữ liệu vào đồ thị. Các mục bị từ chối (\texttt{REJECTED}) sẽ bị loại bỏ ngay lập tức, ngăn chặn việc làm ô nhiễm đồ thị. Các mục hợp lệ (\texttt{APPROVED}) được hợp nhất bình thường. Đối với các mục chờ xem xét (\texttt{PENDING\_REVIEW}), quy trình quản lý dữ liệu chờ duyệt này được đưa vào một vùng đệm lưu trữ gọi là Staging Area (Cơ sở dữ liệu Staging). Tại đây, người quản trị có thể can thiệp thủ công thông qua giao diện quản trị hoặc các API để rà soát chi tiết từng thực thể hoặc mối quan hệ (Staging Review): nếu được phê duyệt (Admin Approved), dữ liệu sẽ được chuyển trạng thái thành \texttt{APPROVED} và hợp nhất vào đồ thị tri thức chính thức; ngược lại nếu bị bác bỏ (Admin Rejected), dữ liệu sẽ chuyển thành \texttt{REJECTED} và bị xóa hoàn toàn khỏi vùng đệm. Việc lưu vết minh bạch từng quyết định (kèm theo đoạn văn bản làm bằng chứng) đảm bảo rằng đồ thị luôn có cơ sở vững chắc và đáng tin cậy. Hình~\ref{fig:flow_shacl_3guard} minh họa quy trình kiểm chứng dựa trên SHACL kết hợp rà soát Staging Review.

\begin{figure}[H]
	\centering
	\begin{tikzpicture}[node distance=1.2cm, auto, scale=0.8, transform shape]
		% Nodes
		\node (in) [startstop, text width=5.5cm] {Thực thể \& Quan hệ trích xuất \\ \footnotesize(theo từng phân đoạn)};
		\node (g1) [decision, below=1cm of in, text width=3cm] {Kiểm tra \texttt{chunk\_id} đã tồn tại?};
		\node (proceed) [smallbox, below left=1.1cm and 1.3cm of g1, text width=3.6cm] {PROCEED \\ \footnotesize(phân đoạn mới, tính bậc cộng dồn)};
		\node (reingest) [smallbox, below right=1.1cm and 1.3cm of g1, text width=3.6cm] {RE\_INGEST \\ \footnotesize(đã tồn tại: nạp lại, tính bậc theo max)};

		\node (g2box) [groupbox, below=3.7cm of g1, minimum width=11cm, minimum height=2.6cm] (g2) {};
		\node [above right=1mm of g2.north west, anchor=south west, font=\small\bfseries] {Kiểm tra cấu trúc \& Đối chiếu bằng chứng};
		\node (shapes) [smallbox, text width=4.6cm] at ($(g2.center)+(-2.7cm,0)$) {SHACL Shapes \\ \footnotesize(name, type, isolated, provenance, weight, foreign-script)};
		\node (ground) [smallbox, text width=4.6cm] at ($(g2.center)+(2.7cm,0)$) {Grounding xác định \\ \footnotesize(exact + fuzzy $\ge$ 0.9, span bằng chứng)};

		\node (dec) [decision, below=1.5cm of g2, text width=3.2cm] {\texttt{shacl\_pass} \& grounded?};

		\node (approved) [startstop, below left=1.4cm and 1.2cm of dec, text width=2.8cm, fill=white] {APPROVED \\ \footnotesize(Hợp nhất Đồ thị)};
		\node (pending) [startstop, below=1.5cm of dec, text width=3cm, fill=white] {PENDING\_REVIEW \\ \footnotesize(Gắn cờ shacl\_review)};
		\node (rejected) [startstop, below right=1.4cm and 1.2cm of dec, text width=2.8cm, fill=white] {REJECTED \\ \footnotesize(Loại bỏ)};
		\node (staging) [smallbox, below=1.2cm of pending, text width=4cm, fill=white] {Staging Area \\ \footnotesize(Giao diện Admin \& API rà soát)};

		% Arrows
		\draw [arrow] (in) -- (g1);
		\draw [arrow] (g1) -| node[anchor=south, pos=0.25, font=\footnotesize] {KHÔNG} (proceed);
		\draw [arrow] (g1) -| node[anchor=south, pos=0.25, font=\footnotesize] {CÓ} (reingest);
		\draw [arrow] (proceed.south) -- ([xshift=-2.7cm]g2.north);
		\draw [arrow] (reingest.south) -- ([xshift=2.7cm]g2.north);
		\draw [arrow] (g2) -- (dec);
		\draw [arrow] (dec) -| node[anchor=south, pos=0.25, font=\footnotesize] {Đủ bằng chứng} (approved);
		\draw [arrow] (dec) -- node[anchor=west, font=\footnotesize] {Thiếu BC, LLM$\ge$0.8} (pending);
		\draw [arrow] (dec) -| node[anchor=south, pos=0.25, font=\footnotesize] {Lỗi cấu trúc} (rejected);
		\draw [arrow] (pending) -- (staging);
		\draw [arrow, dashed, rounded corners] (staging.west) -| node[anchor=north, pos=0.25, font=\footnotesize] {Admin duyệt} (approved.south);
		\draw [arrow, dashed, rounded corners] (staging.east) -| node[anchor=north, pos=0.25, font=\footnotesize] {Admin bác bỏ} (rejected.south);
	\end{tikzpicture}
	\caption{Quy trình kiểm chứng tri thức SHACL và đối chiếu bằng chứng xác định}
	\label{fig:flow_shacl_3guard}
\end{figure}

% ------------------------------------------------------------------
\subsection{Xử lý thực thể trùng lặp trong văn bản học thuật}
% ------------------------------------------------------------------

% ------------------------------------------------------------------
\label{subsubsec:entity-resolution-solution}
% ------------------------------------------------------------------

Để xử lý vấn đề thực thể trùng lặp đảm bảo truy xuất đầy đủ thông tin, đề tài đề xuất một module phân giải thực thể
(\textit{entity resolution}) được tích hợp vào hai giai đoạn,
hoạt động trước khi bất kỳ nút nào được ghi vào đồ thị. Module này
được thiết kế theo kiến trúc hai giai đoạn xử lý nối tiếp:
Candidate Retrieval (Truy hồi ứng viên) và LLM
Decision(LLM đưa ra quyết định).

Cách tiếp cận này được xây dựng trên nền tảng của hai công trình nghiên
cứu trong lĩnh vực căn chỉnh thực thể (\textit{entity alignment}) giữa
các knowledge graph: HELEA~\cite{helea2024} chứng minh hiệu quả của
    pipeline hai giai đoạn: đầu tiên truy hồi ứng viên bằng embedding
    similarity, sau đó dùng mô hình ngôn ngữ lớn để xếp hạng lại
    (\textit{reranking}). Cơ chế này vừa đảm bảo không bỏ sót ứng viên
    tiềm năng, vừa tận dụng khả năng suy luận ngữ nghĩa sâu của LLM. Kế tiếp là EasyEA~\cite{easyea2024} chứng minh rằng mô hình
    ngôn ngữ lớn, khbi được cung cấp đủ ngữ cảnh, có khả năng đưa ra
    quyết định căn chỉnh thực thể chính xác theo phương thức
    \textit{zero-shot} kbhông cần dữ liệu huấn luyện nhãn, phù hợp
    với bối cảnh corpus tiếng Việt học thuật hiện tại. Trong Hình~\ref{fig:entity-resolution-flow} Trình bày kiến trúc tổng thể của
nhất đồ thị.

% ---- TIKZ DIAGRAM ----
\begin{figure}[H]
\centering
\caption{Kiến trúc hai giai đoạn của module phân giải thực thể.
Giai đoạn Candidate Retrieval sử dụng embedding similarity để truy hồi
các thực thể ngữ nghĩa tương đồng từ vector store. Giai đoạn LLM
Decision chỉ được kích hoạt khi có ứng viên vượt ngưỡng $\tau$,
giúp tiết kiệm chi phí gọi LLM cho những thực thể hoàn toàn mới.}
\label{fig:entity-resolution-flow}
\begin{tikzpicture}[
    node distance=0.55cm and 1.5cm,
    every node/.style={font=\small},
    %
    inputbox/.style={
        rectangle, rounded corners=4pt,
        draw=black, fill=gray!10,
        text width=9cm, align=center,
        minimum height=0.9cm, inner sep=6pt
    },
    stagebox/.style={
        rectangle, rounded corners=5pt,
        draw=black, fill=white,
        text width=9cm, align=left,
        inner sep=10pt
    },
    decisionbox/.style={
        diamond, draw=black, fill=white,
        text width=2.6cm, align=center,
        aspect=2.8, inner sep=2pt
    },
    skipbox/.style={
        rectangle, rounded corners=3pt,
        draw=black, fill=white, dashed,
        text width=3.2cm, align=center,
        minimum height=0.8cm, inner sep=5pt
    },
    outputbox/.style={
        rectangle, rounded corners=4pt,
        draw=black, fill=gray!15,
        text width=9cm, align=center,
        minimum height=0.9cm, inner sep=6pt
    },
    arrow/.style={-Stealth, thick, color=black},
    dasharrow/.style={-Stealth, thick, dashed, color=black}
]

%% INPUT
\node[inputbox] (input) {%
    Thực thể đầu vào từ giai đoạn trích xuất\\
    \small\{tên thực thể, mô tả ngữ nghĩa, loại thực thể\}
};

%% STAGE 1 — Candidate Retrieval
\node[stagebox, below=0.55cm of input] (s1) {%
    Giai đoạn 1 : Candidate Retrieval\\[6pt]
    Mỗi thực thể đầu vào được biểu diễn dưới dạng một vector embedding
    bằng cách kết hợp tên và mô tả ngữ nghĩa. Vector này được dùng
    để truy vấn vector store nhằm tìm tối đa $K$ thực thể đã tồn
    tại có ngữ nghĩa gần nhất. Chỉ những ứng viên có cosine
    similarity vượt ngưỡng $\tau$ mới được chuyển sang giai đoạn
    tiếp theo.
};

%% DECISION
\node[decisionbox, below=0.55cm of s1] (dec) {
    \footnotesize Tìm được\\ứng viên\\tiềm năng?
};

%% SKIP
\node[skipbox, right=1.4cm of dec] (skip) {%
    \centering\small Thực thể mới\\hoàn toàn\\$\rightarrow$ Tạo nút mới
};

%% STAGE 2 — LLM Decision
\node[stagebox, below=0.55cm of dec] (s2) {%
    Giai đoạn 2 : LLM Decision\\[6pt]
    Mô hình ngôn ngữ lớn nhận đồng thời thông tin của thực thể
    mục tiêu và danh sách ứng viên. Dựa trên phân tích ngữ nghĩa
    tổng hợp, bao gồm so sánh tên gọi, mô tả, và các quy tắc
    phân biệt tường minh, LLM đưa ra phán quyết cuối cùng: hai
    thực thể có đồng nhất hay không, và nếu có, thực thể nào là
    dạng chuẩn (\textit{canonical form}).
};

%% OUTPUT
\node[outputbox, below=0.55cm of s2] (output) {%
    Kết quả phân giải:
    ánh xạ đồng nhất hoặc tạo nút mới độc lập
};

%% ARROWS
\draw[arrow] (input)  -- (s1);
\draw[arrow] (s1)     -- (dec);
\draw[arrow] (dec)    -- node[left, font=\footnotesize]{Có} (s2);
\draw[dasharrow] (dec) -- node[above, font=\footnotesize]{Không} (skip);
\draw[arrow] (s2)     -- (output);
\draw[arrow] (skip.south) |- (output.east);

\end{tikzpicture}
\end{figure}

Giai đoạn 1: Candidate Retrieval.
Mỗi thực thể cần phân giải được biểu diễn dưới dạng một vector embedding bằng cách kết hợp cả tên lẫn
mô tả của thực thể đó thành một chuỗi đầu vào thống nhất. Cách biểu
diễn này cho phép vector nắm bắt đồng thời hai chiều thông tin: chiều
định danh (tên gọi) và chiều ngữ nghĩa (nội dung, thuộc tính). Sau đó,
hệ thống thực hiện tìm kiếm láng giềng gần nhất trong vector store để truy hồi tối đa $K$ thực thể đã
tồn tại trong knowledge graph có vector gần nhất với thực thể đang xét.
Ngưỡng tương đồng cosine $\tau$ đóng vai trò bộ lọc: chỉ những ứng
viên có $\cos(\mathbf{v}_1, \mathbf{v}_2) \geq \tau$ mới được coi là
\textit{ứng viên tiềm năng} và đưa vào giai đoạn tiếp theo. Cơ chế này tương đồng
với giai đoạn truy hồi ứng viên trong kiến trúc HELEA~\cite{helea2024},
đảm bảo không bỏ sót các cặp thực thể có ngữ nghĩa gần nhau.

Giai đoạn 2: LLM Decision.
Khi giai đoạn Candidate Retrieval xác định được ít nhất một ứng viên
tiềm năng, mô hình ngôn ngữ lớn được kích hoạt để đưa ra phán quyết
cuối cùng. Không giống với các phương pháp truyền thống chỉ dựa vào
điểm số cosine similarity, LLM có khả năng suy luận ở mức ngữ nghĩa
sâu hơn, cũng như phân biệt các thực thể bề ngoài tương tự
nhưng thực chất khác biệt. Chẳng hạn, mô hình xác định
\textit{``Nhập môn Lập trình''} và \textit{``Lập trình Hướng đối
tượng''} là hai học phần khác nhau mặc dù đều thuộc nhóm
lập trình, trong khi \textit{``HCMUS''} và \textit{``Trường Đại học
Khoa học Tự nhiên''} là cùng một tổ chức.
% ------------------------------------------------------------------
\subsubsection{Thiết kế quan hệ SAME\_AS: Liên kết thay vì xóa bỏ}
\label{subsubsec:sameas-design}
% ------------------------------------------------------------------

Khi mô hình ngôn ngữ lớn xác nhận hai thực thể là đồng nhất, hệ thống có 2 lựa chọn: \textit{hợp nhất vật lý}
hai nút thành một (xóa một nút, gộp toàn bộ dữ liệu vào nút còn lại),
hay \textit{tạo quan hệ liên kết} giữa hai thực thể này và giữ nguyên cả hai? Đề
tài lựa chọn giải pháp thứ hai thông qua quan hệ SAME\_AS.

Quyết định này không phải là sự thỏa hiệp mà xuất phát từ ba lý do cơ
bản. \textit{Thứ nhất}, cả HELEA lẫn EasyEA đều tiếp cận entity
alignment theo hướng tạo ra \textit{tập hợp cặp đồng nhất} giữa hai
knowledge graph độc lập không ai trong số họ đề xuất xóa bỏ hoặc
hợp nhất vật lý các nút. Quan hệ SAME\_AS là biểu diễn ngữ nghĩa
chuẩn của kết quả alignment, tương đương với thuộc tính
owl:sameAs trong các chuẩn ngữ nghĩa web (\textit{Semantic
Web})~\cite{owlsemweb2004}. \textit{Thứ hai}, hợp nhất vật lý là
thao tác không thể hoàn tác: nếu mô hình ngôn ngữ lớn đưa ra quyết
định nhầm, dữ liệu của nút bị xóa không thể phục hồi.
Ngược lại, một quan hệ SAME\_AS có thể được xóa bỏ đơn giản nếu phát
hiện sai sót, trong khi toàn bộ thông tin gốc vẫn được bảo toàn.
\textit{Thứ ba}, thiết kế SAME\_AS bảo toàn \textit{nguồn gốc ngữ
nghĩa} (\textit{provenance}): mỗi nút vẫn duy trì liên kết đến đúng
các tài liệu đã sử dụng hình thức tên đó, phục vụ cho việc truy vết
và kiểm chứng thông tin sau này.

Về mặt hành vi trong knowledge graph, thiết kế này cho phép hệ thống
truy vấn tổng hợp thông tin từ cả hai nút thông qua phép duyệt cạnh
SAME\_AS. Ví dụ, khi người dùng hỏi về \textit{``HCMUS''}, hệ thống
không chỉ truy xuất thông tin gắn với nút \textit{``HCMUS''} mà còn
có thể traverse sang nút \textit{``Trường Đại học Khoa học Tự nhiên''}
để bổ sung thêm tri thức liên quan, cho ra câu trả lời toàn diện hơn.

% ------------------------------------------------------------------
\subsubsection{Kết quả trong knowledge graph}
\label{subsubsec:kg-result}
% ------------------------------------------------------------------

Minh họa cấu trúc của knowledge graph sau
khi áp dụng module phân giải thực thể với dữ liệu từ corpus học thuật.

\begin{figure}[H]
\centering
\resizebox{\linewidth}{!}{%
\begin{tikzpicture}[
    every node/.style={font=\normalsize},
    canonical/.style={
        rectangle, rounded corners=6pt,
        draw=black, fill=gray!25, line width=1.2pt,
        minimum width=5.0cm, minimum height=1.0cm,
        align=center, inner sep=7pt, font=\normalsize\bfseries
    },
    alias/.style={
        rectangle, rounded corners=5pt,
        draw=black!70, fill=white,
        minimum width=3.0cm, minimum height=0.9cm,
        align=center, inner sep=6pt
    },
    concept/.style={
        rectangle, rounded corners=4pt,
        draw=black!70, fill=gray!10,
        minimum width=3.2cm, minimum height=0.85cm,
        align=center, inner sep=5pt
    },
    sameas/.style={-Stealth, thick, dashed, draw=black},
    rel/.style={-Stealth, thick, draw=black},
    lbl/.style={font=\small\itshape, fill=white, inner sep=2pt},
    slbl/.style={font=\small\bfseries, color=black, fill=white, inner sep=2pt},
]

% ── Hàng trên: concept nodes ───────────────────────────────────────────
\node[concept] (ai)     at ( 0.0,  3.2) {Trí tuệ nhân tạo};
\node[concept] (dhqg)   at ( 5.0,  3.2) {ĐHQG-HCM};
\node[concept] (thuat)  at (11.0,  3.2) {Giải thuật};
\node[concept] (pttkht) at (16.0,  3.2) {PT\&TK Hệ thống};

% ── Hàng giữa: canonical nodes ─────────────────────────────────────────
\node[canonical] (school) at ( 2.5,  0.0) {Trường ĐH Khoa học Tự nhiên};
\node[canonical] (csdl_c) at (13.5,  0.0) {Cơ sở dữ liệu};

% ── Hàng dưới: alias nodes ─────────────────────────────────────────────
\node[alias] (khtn)   at (-1.0, -3.0) {Trường ĐH KHTN};
\node[alias] (hcmus)  at ( 2.8, -3.0) {HCMUS};
\node[alias] (uos)    at ( 6.5, -3.0) {University of Science};
\node[alias] (csdl_a) at (11.5, -3.0) {CSDL};
\node[alias] (db)     at (15.5, -3.0) {Database};

% ── SAME_AS edges ──────────────────────────────────────────────────────
\draw[sameas] (khtn.north)
    to[bend left=8]
    node[slbl, midway, left=3pt]{SAME\_AS}
    (school.south west);

\draw[sameas] (hcmus.north)
    -- node[slbl, midway, right=3pt]{SAME\_AS}
    (school.south);

\draw[sameas] (uos.north)
    to[bend right=15]
    node[slbl, midway, right=4pt]{SAME\_AS}
    (school.south east);

\draw[sameas] (csdl_a.north)
    -- node[slbl, midway, left=3pt]{SAME\_AS}
    (csdl_c.south west);

\draw[sameas] (db.north)
    to[bend right=8]
    node[slbl, midway, right=3pt]{SAME\_AS}
    (csdl_c.south east);

% ── Quan hệ ngữ nghĩa ──────────────────────────────────────────────────
\draw[rel] (school.north west)
    to[bend right=15]
    node[lbl, midway, left=3pt]{giảng dạy}
    (ai.south);

\draw[rel] (school.north east)
    to[bend left=15]
    node[lbl, midway, right=3pt]{thuộc}
    (dhqg.south);

\draw[rel] (csdl_c.north west)
    to[bend right=15]
    node[lbl, midway, left=3pt]{tiên quyết của}
    (thuat.south);

\draw[rel] (csdl_c.north east)
    to[bend left=15]
    node[lbl, midway, right=3pt]{tiên quyết của}
    (pttkht.south);

\end{tikzpicture}
}% end resizebox
\caption{Ví dụ cấu trúc knowledge graph sau khi áp dụng phân giải thực
thể trên corpus học thuật HCMUS. Các nút viền đậm, màu xám đậm là thực thể dạng
chuẩn (\textit{canonical}); các nút màu trắng là hình thức tên thay
thế (\textit{alias}) được liên kết về nút chuẩn qua quan hệ SAME\_AS
(đường nét đứt). Quan hệ ngữ nghĩa thông thường (mũi tên nét liền)
chỉ xuất phát từ nút chuẩn, đảm bảo tri thức không bị phân tán.}
\label{fig:sameas-graph}
\end{figure}

% ------------------------------------------------------------------
\subsubsection{Đồng bộ hóa với vector store}
% ------------------------------------------------------------------

Song song với việc xây dựng đồ thị quan hệ, hệ thống duy trì một cơ sở
dữ liệu vector (\textit{vector store}) được đồng bộ liên tục. Sau mỗi
thao tác hợp nhất, toàn bộ thực thể và quan hệ - với mô tả đã được
tổng hợp - đều được cập nhật vào vector store. Kiến trúc lưu trữ kép
này (\textit{dual storage}) là nền tảng cho hai chế độ truy vấn của
LightRAG: đồ thị tri thức phục vụ duyệt cấu trúc quan hệ tường minh,
trong khi vector store cho phép tìm kiếm ngữ nghĩa linh hoạt khi câu
hỏi của người dùng không khớp chính xác với bất kỳ tên thực thể nào
đã được lập chỉ mục.

% ------------------------------------------------------------------
\subsubsection{So sánh với công trình liên quan}
\label{subsubsec:comparison-related}
% ------------------------------------------------------------------

Bảng~\ref{tab:er-comparison} tổng hợp mức độ tương đồng về mặt thiết
kế giữa module phân giải thực thể của đề tài và hai công trình nghiên
cứu tham chiếu.

\begin{table}[H]
\centering
\caption{So sánh các thành phần thiết kế với công trình liên quan}
\label{tab:er-comparison}

\begin{tabular}{p{4.5cm} >{\centering\arraybackslash}p{2cm} >{\centering\arraybackslash}p{2cm} p{6cm}}
\toprule
Thành phần &
HELEA &
EasyEA &
Phương án đề tài \\
\midrule

Embedding kết hợp tên \& mô tả
& \ding{51} & \ding{51}
& Áp dụng. Biểu diễn thực thể qua vector kết hợp hai chiều ngữ nghĩa. \\
\addlinespace

Truy hồi Top-$K$ qua vector similarity
& \ding{51} & \ding{51}
& Áp dụng với lọc ngưỡng cosine $\tau$. \\
\addlinespace

LLM ra quyết định cuối cùng
& \ding{51} & \ding{51}
& Áp dụng theo zero-shot. \\
\addlinespace

Quan hệ SAME\_AS (thay vì merge vật lý)
& \ding{51} & \ding{51}
& Áp dụng. Bảo toàn cả hai nút và provenance. \\
\addlinespace

Ngữ cảnh triple láng giềng trong prompt
& $-$ & \ding{51}
& Chưa tích hợp. Xác định là hướng phát triển tiếp theo. \\
\addlinespace

Fine-tuning encoder với hard negatives
& \ding{51} & $-$
& Không áp dụng (chọn hướng zero-shot). \\

\bottomrule
\end{tabular}
\end{table}

Hai điểm khác biệt chính so với các công trình tham chiếu cần được
làm rõ. Đối với việc không fine-tuning encoder, đây là quyết định có
chủ ý: corpus tiếng Việt trong lĩnh vực học thuật gần như không có tập
dữ liệu căn chỉnh thực thể được gán nhãn, khiến việc huấn luyện không
khả thi. Phương án zero-shot mà EasyEA kiểm chứng giải quyết đúng hạn
chế này. Đối với ngữ cảnh triple láng giềng, EasyEA chỉ ra rằng việc
bổ sung các bộ ba quan hệ liên quan của ứng viên vào prompt giúp LLM
phân biệt tốt hơn các trường hợp thực thể mơ hồ - ví dụ, một thực
thể là \textit{học phần} và một thực thể là \textit{phần mềm} có thể
cùng mang tên ``Cơ sở dữ liệu'' nhưng có bộ quan hệ hoàn toàn khác
nhau trong knowledge graph. Việc tích hợp thành phần này được xác định
là hướng phát triển ưu tiên trong giai đoạn tiếp theo.
\subsection{Quản lý và duy trì bền vững hệ thống}

\indent LightRAG thiết kế hệ thống quản trị với triết lý nhất quán tăng dần (incremental consistency): thay vì xây dựng lại toàn bộ đồ thị mỗi khi có tài liệu mới, hệ thống chỉ xử lý các đoạn văn mới hoặc đã thay đổi, sau đó hợp nhất kết quả vào đồ thị hiện có. Cơ chế làm sạch tự động (auto-cleaning) thường xuyên kiểm tra tính nhất quán giữa đồ thị tri thức và nội dung tài liệu gốc: nếu một tài liệu bị xóa khỏi kho dữ liệu, tất cả các nút và cạnh có nguồn gốc duy nhất từ tài liệu đó sẽ tự động bị loại bỏ, trong khi các nút được hỗ trợ bởi nhiều nguồn sẽ chỉ được cập nhật lại mô tả.

\indent Người quản trị có thể can thiệp thủ công thông qua các API quản trị: hiệu chỉnh nhãn thực thể (entity label editing), cập nhật mô tả hoặc hợp nhất thủ công hai thực thể bị nhận diện nhầm là khác nhau. Sau bất kỳ thay đổi nào, hệ thống kích hoạt quy trình đồng bộ hóa tự động: véc-tơ nhúng của các thực thể bị ảnh hưởng được tính toán lại và cập nhật trong cơ sở dữ liệu vector, đảm bảo sự nhất quán tuyệt đối giữa hai lớp lưu trữ.

\indent Cuối cùng, LightRAG hỗ trợ tùy chỉnh lời nhắc trích xuất theo lĩnh vực (domain-specific prompt tuning): người dùng có thể cung cấp danh sách các kiểu thực thể ưu tiên và các mẫu quan hệ đặc thù, từ đó định hướng mô hình tập trung vào tri thức có giá trị nhất trong từng miền chuyên môn cụ thể.

\section{Mô hình hóa trạng thái hệ thống (Agent State)}
Phân hệ trợ lý ảo được xây dựng trên nền tảng LangGraph và thiết kế dưới dạng một đồ thị trạng thái có hướng (state graph), trong đó mỗi nút đại diện cho một tác vụ tính toán và các cạnh đại diện cho luồng điều khiển. Trạng thái dùng chung của hệ thống được mô hình hóa bằng kiểu \texttt{AgentState} và có thể biểu diễn bởi tập hợp $S = \{Q, U, H, A, D, K, M, G, F\}$. Trong đó, $Q$ (Query) là câu hỏi đầu vào, $U$ (User ID) là mã định danh người dùng phục vụ cá nhân hóa, và $H$ (Chat History) là lịch sử hội thoại ngắn hạn đã được tóm tắt. Thành phần $A$ (Analysis) lưu kết quả phân tích ý định, $D$ (Decision) là quyết định điều phối với ba giá trị \texttt{direct\_answer}, \texttt{retrieve} hoặc \texttt{reject}, còn $K$ (Keywords) là tập từ khóa truy xuất được tách thành hai mức cao và thấp. Cuối cùng, $M$ (Long-term Memory) là hồ sơ người dùng dài hạn, $G$ (Generation Request) là gói ngữ cảnh đầy đủ để sinh câu trả lời, và $F$ (Final Response) là câu trả lời cuối cùng trả về cho người dùng. So với các kiến trúc có vòng lặp tự đánh giá, hệ thống đề xuất chủ trương giữ luồng điều phối tuyến tính và tất định để tối ưu độ trễ phản hồi, dồn chất lượng vào khâu phân tích đầu vào và truy xuất tri thức thay vì hậu kiểm nhiều vòng.

\subsection{Kiến trúc các tác vụ (Agents)}

Hệ thống bao gồm năm thành phần xử lý chuyên biệt, được hiện thực dưới dạng các hàm trạng thái (stateless functions) và phối hợp theo một đồ thị có hướng do LangGraph điều phối, thay vì các tác nhân tự hành lập kế hoạch nhiều lượt.

Thứ nhất, Analyze Agent (\texttt{analyze\_agent}) là thành phần trung tâm của khâu định tuyến. Bằng một lần gọi mô hình ngôn ngữ lớn với định dạng đầu ra JSON, agent này phân loại ý định người dùng, đưa ra quyết định điều phối (\texttt{direct\_answer}, \texttt{retrieve} hoặc \texttt{reject}), chọn chế độ truy vấn LightRAG, viết lại câu hỏi thành dạng độc lập (standalone query), và trích xuất hai tập từ khóa: từ khóa mức thấp (cụ thể, như mã môn học, tên giảng viên) và mức cao (khái quát, như chính sách, so sánh). Agent cũng đảm nhiệm vai trò kiểm duyệt: khi phát hiện nội dung không phù hợp, quyết định bị ép sang \texttt{reject}.

Thứ hai, Generator Agent (\texttt{generator}) chịu trách nhiệm dựng ngữ cảnh sinh chứ không trực tiếp gọi mô hình. Với nhánh truy xuất, persona của trợ lý được đặt vào trường \texttt{user\_prompt} để LightRAG vẫn giữ được khuôn mẫu nội bộ và chèn ngữ cảnh tri thức \texttt{context\_data}; với nhánh trả lời trực tiếp, agent dựng một \texttt{system\_prompt} đầy đủ; còn với nhánh từ chối, agent trả về một thông điệp từ chối lịch sự đã soạn sẵn.

Thứ ba, Summarizer Agent (\texttt{summarizer}) nén lịch sử hội thoại khi vượt quá bốn lượt thành một đến hai câu súc tích, nhằm giảm số token mà Analyze Agent phải xử lý và duy trì ngữ cảnh cốt lõi qua các phiên dài.

Thứ tư, Long-term Memory Agent (\texttt{longterm}) quản lý hồ sơ người dùng bền vững trên MongoDB. Thay vì tóm tắt lại toàn bộ hồ sơ, agent chỉ trích xuất các sự kiện ổn định (chuyên ngành, sở thích trả lời, định hướng học tập dài hạn) dưới dạng JSON rồi hợp nhất các thông tin này vào hồ sơ hiện có bằng mã Python có khử trùng lặp, tránh hiện tượng trôi dạt nội dung tóm tắt.

Cuối cùng, LLM Factory (\texttt{llm}) là lớp trừu tượng hóa nhà cung cấp mô hình, cho phép chuyển đổi linh hoạt giữa OpenAI và Ollama (mô hình cục bộ) qua biến môi trường và tự chuẩn hóa các tham số khác biệt giữa hai nền tảng.

\subsection{Quy trình điều phối và Luồng dữ liệu}

Đồ thị điều phối được biên dịch cùng một bộ lưu điểm kiểm tra (checkpointer) trên MongoDB và gồm bốn nút thực thi tuần tự: \texttt{analyzer} $\rightarrow$ \texttt{prepare\_context} $\rightarrow$ \texttt{generate} $\rightarrow$ (có điều kiện) \texttt{update\_long\_term} $\rightarrow$ \texttt{END}. Quy trình thực thi được mô tả qua các bước logic sau:

Bước 1: Phân tích và Định tuyến.
Khi nhận câu hỏi $Q$, nút \texttt{analyzer} gọi \textit{Analyze Agent} (sau khi đã tóm tắt lịch sử qua \textit{Summarizer}) để xác định ý định, quyết định điều phối, chế độ truy vấn và các từ khóa. Kết quả phân tích được chuẩn hóa: nếu nội dung không phù hợp thì quyết định bị ép sang \texttt{reject}; nếu giá trị không hợp lệ thì lùi an toàn về \texttt{retrieve}.

Bước 2: Chuẩn bị ngữ cảnh.
Nút \texttt{prepare\_context} không gọi mô hình ngôn ngữ. Nút này chỉ nạp bộ nhớ dài hạn $M$ từ MongoDB khi khâu phân tích xác định là cần thiết, rồi dựng gói yêu cầu sinh $G$ (\texttt{generation\_request}) chứa câu hỏi, chế độ truy vấn, persona và các từ khóa. Việc tách riêng bước này giúp luồng phát trực tuyến (streaming) có thể dừng tại đây trước khi sinh.

Bước 3: Sinh câu trả lời.
Nút \texttt{generate} phân nhánh theo quyết định điều phối: với \texttt{retrieve}, hệ thống gọi LightRAG kèm các từ khóa đã trích sẵn để truy xuất tri thức từ đồ thị; với \texttt{direct\_answer}, hệ thống gọi trực tiếp mô hình ngôn ngữ mà không truy xuất (phù hợp cho chào hỏi, trò chuyện hoặc lời khuyên học tập chung); với \texttt{reject}, hệ thống trả về thông điệp từ chối đã soạn sẵn.

Bước 4: Cập nhật bộ nhớ dài hạn (có điều kiện).
Khác với cập nhật theo chu kỳ cố định, hệ thống chỉ kích hoạt nút \texttt{update\_long\_term} khi khâu phân tích đánh dấu cờ \texttt{should\_update\_long\_term} - tức khi người dùng cung cấp một sự kiện ổn định đáng ghi nhớ. Khi đó, \textit{Long-term Memory Agent} trích xuất và hợp nhất thông tin mới vào hồ sơ; ngược lại, đồ thị kết thúc ngay sau bước sinh để giảm độ trễ.

Quy trình điều phối tổng thể và luồng dữ liệu giữa các phân hệ được mô tả trực quan trong Hình~\ref{fig:flow_query_lightrag}.

\begin{figure}[H]
	\centering
	\resizebox{\textwidth}{!}{
		\begin{tikzpicture}[
			>={Stealth[scale=1.2]},
			base/.style={align=center, font=\sffamily\small, inner sep=0.3cm, thick},
			inputQuery/.style={base, rectangle, rounded corners=5pt, fill=black!80, text=white, draw=black, minimum width=3.5cm},
			outputResp/.style={base, rectangle, rounded corners=5pt, fill=black!80, text=white, draw=black, minimum width=4cm},
			retrieveBox/.style={base, rectangle, rounded corners=5pt, fill=white, draw=black, minimum width=3.5cm},
			ragSystemBox/.style={base, rectangle, rounded corners=5pt, fill=white, draw=black, minimum width=5.5cm, minimum height=1.5cm},
			generateBox/.style={base, rectangle, rounded corners=5pt, fill=white, draw=black, minimum width=4.5cm},
			storageBox/.style={base, cylinder, shape border rotate=90, aspect=0.2, fill=gray!30, text=black, draw=black, minimum height=1.3cm, minimum width=4cm},
			memoryBox/.style={base, cylinder, shape border rotate=90, aspect=0.2, fill=white, draw=black, minimum height=1.3cm, text width=4cm},
			llmBox/.style={base, rectangle, rounded corners=5pt, fill=white, draw=black, minimum width=4.5cm},
			decisionBox/.style={base, diamond, aspect=2.5, fill=white, draw=black, inner sep=2pt, text width=3cm},
			groupLabel/.style={font=\bfseries\sffamily\large, right, fill=white, inner sep=3pt, rounded corners=3pt},
			labelTag/.style={font=\scriptsize, fill=white, inner sep=1.5pt, rounded corners=2pt} % Nhãn đè lên vạch không bị rối
		]

			% ==========================================
			% LAYER 0: BACKGROUND GROUPS 
			% ==========================================
			\begin{scope}[on background layer]
				% 0. Pre-Process Group
				\filldraw[fill=gray!3, draw=black!40, dashed, rounded corners=15pt, very thick] (-15, 6.8) rectangle (16, 0.5);
				\node[groupLabel, text=black] at (-14.5, 6.8) {0. Phân tích \& Định tuyến (Analyzer Agent)};
				
				% 1. Retrieve Group 
				\filldraw[fill=gray!3, draw=black!40, dashed, rounded corners=15pt, very thick] (-15, 0) rectangle (16, -4.5);
				\node[groupLabel, text=black] at (-14.5, 0) {1. Truy xuất Ngữ cảnh (LightRAG System)};
				
				% 2. Generation Group
				\filldraw[fill=gray!3, draw=black!40, dashed, rounded corners=15pt, very thick] (-15, -5) rectangle (16, -7.5);
				\node[groupLabel, text=black] at (-14.5, -5) {2. Tạo sinh câu trả lời (Generator Agent)};
				
				% 3. Memory Group
				\filldraw[fill=gray!3, draw=black!40, dashed, rounded corners=15pt, very thick] (-1, -8) rectangle (16, -11);
				\node[groupLabel, text=black] at (-0.5, -8) {3. Cập nhật Bộ nhớ (Memory Agent)};
			\end{scope}

			% ==========================================
			% 0. PRE-PROCESS \& ROUTING PHASE 
			% ==========================================
			\node[inputQuery] (Query) at (-12, 1.5) {Query\\(Câu hỏi)};
			\node[memoryBox, text width=3.5cm] (ChatHistory) at (-6, 4.5) {Chat History\\(Lịch sử ngắn hạn)};
			\node[llmBox, text width=6cm] (AnalyzerLLM) at (-6, 1.5) {Analyzer LLM\\1. Làm rõ câu hỏi\\2. Phân loại ý định\\3. Trích xuất Keywords};
			\node[decisionBox] (RoutingDecision) at (0, 1.5) {Decision\\(Routing)};
			\node[decisionBox, text width=4cm] (ModeSelection) at (7, 1.5) {RAG Mode Selection\\\scriptsize(local, global, hybrid, mix)\\\scriptsize*Fallback mặc định: mix};

			% ==========================================
			% 1. RETRIEVAL PHASE (LightRAG call simplified)
			% ==========================================
			\node[ragSystemBox, text width=6.5cm] (LightRAGSystem) at (0, -2.2) {LightRAG Core System\\- Tìm kiếm Vector (Entities/Relations)\\- Duyệt đồ thị quan hệ\\- Lọc nhiễu và Trích xuất ngữ cảnh};

			% ==========================================
			% 2. GENERATION PHASE
			% ==========================================
			\node[generateBox] (CombCtx) at (-5.5, -6.2) {Combined Context\\(Ngữ cảnh kết hợp)};
			\node[generateBox, text width=5.5cm] (GenPrompt) at (1.5, -6.2) {Build Generation Request\\(Inject Memory \& Context)};
			\node[llmBox] (GeneratorLLM) at (7.5, -6.2) {Generator LLM};
			\node[outputResp] (Response) at (13, -6.2) {Response\\(Câu trả lời)};

			% ==========================================
			% 3. MEMORY UPDATE 
			% ==========================================
			\node[llmBox, text width=3.5cm] (MemUpdateLLM) at (13, -9.5) {Update Profile LLM\\(Phân tích)};
			\node[memoryBox, text width=3.5cm] (LongTermMem) at (1.5, -9.5) {Long-Term Memory\\(Hồ sơ user)};

			% ==========================================
			% WIRING \& ROUTING
			% ==========================================
			
			% --- Phase 0 ---
			\draw[->, thick] (Query) -- (AnalyzerLLM);
			\draw[->, thick, dashed, draw=black] (ChatHistory.south) -- node[labelTag, right] {Tóm tắt} (AnalyzerLLM.north);
			\draw[->, thick] (AnalyzerLLM) -- (RoutingDecision);
			\draw[->, thick] (RoutingDecision) -- node[labelTag, above] {decision == 'retrieve'} (ModeSelection);

			% Routing Mode -> LightRAG Call
			\draw[->, thick, rounded corners] (ModeSelection.south) |- (0, -0.5) -- node[labelTag, right, pos=0.85] {Mode: local/global/mix} (LightRAGSystem.north);

			% --- Bypass Lines ---
			% 1. Lịch sử chat -> Prompt (Vành đai Trái)
			\draw[->, thick, dashed, draw=black, rounded corners] 
				(ChatHistory.west) -- (-14.3, 4.5) -- (-14.3, -4.8) -- (0.5, -4.8) -- ([xshift=-1cm]GenPrompt.north);
			\node[labelTag, text=black] at (-11.5, 4.5) {Recent Messages};

			% 2. Direct Answer (Vành đai trên -> Vành đai Phải trong)
			\draw[->, very thick, draw=black, rounded corners] 
				(RoutingDecision.north) -- (0, 3.2) -- (14.5, 3.2) -- (14.5, -4.8) -- (2.5, -4.8) -- ([xshift=1cm]GenPrompt.north);
			\node[labelTag, text=black] at (11, 3.2) {decision == 'direct\_answer'};

			% 3. Reject (Vành đai Phải ngoài cùng)
			\draw[->, very thick, draw=black, rounded corners] 
				(RoutingDecision.north) -- (0, 3.9) -- (15.5, 3.9) -- (15.5, -6.2) -- (Response.east);
			\node[labelTag, text=black] at (11, 3.9) {decision == 'reject'};

			% --- Retrieval Output to Context ---
			\draw[->, thick, rounded corners] (LightRAGSystem.south) -- (0, -4.8) -- (-5.5, -4.8) -- (CombCtx.north);

			% --- Phase 2: Generation ---
			\draw[->, thick] (CombCtx.east) -- (GenPrompt.west);
			\draw[->, thick] (GenPrompt.east) -- (GeneratorLLM.west);
			\draw[->, thick] (GeneratorLLM.east) -- (Response.west);

			% --- Phase 3: Memory Updates ---
			\draw[->, thick, dashed, draw=black] (Response.south) -- node[labelTag, right] {Trigger Update} (MemUpdateLLM.north);
			\draw[->, thick, dashed, draw=black] (MemUpdateLLM.west) -- node[labelTag, above] {Save Summary} (LongTermMem.east);
			\draw[->, thick, dashed, draw=black] (LongTermMem.north) -- node[labelTag, right] {Load Profile} (GenPrompt.south);

		\end{tikzpicture}
	}
	\caption{Sơ đồ kiến trúc điều phối Agent tổng thể của trợ lý ảo}
	\label{fig:flow_query_lightrag}
\end{figure}

\subsection{Cơ chế truy xuất tri thức (LightRAG Integration)}
Khác với RAG truyền thống dựa thuần túy trên độ tương đồng vector, hệ thống sử dụng LightRAG để truy xuất trên cấu trúc đồ thị tri thức (Knowledge Graph). Điểm tối ưu quan trọng là các từ khóa đã được Analyze Agent trích xuất sẵn ở mức cao và mức thấp được truyền thẳng vào LightRAG:
$$Q_{RAG} = Q_{standalone} + K_{high} + K_{low}$$
Nhờ vậy, LightRAG bỏ qua được bước tự trích xuất từ khóa nội bộ (vốn cần thêm một lần gọi mô hình), giúp giảm độ trễ mà vẫn định hướng truy xuất chính xác. Tùy theo độ khái quát và mục đích tra cứu của câu hỏi, hệ thống hỗ trợ bốn chế độ truy vấn chính bao gồm:

\indent Chế độ truy vấn cơ bản (Naive Query): Đây là phương pháp tìm kiếm truyền thống dựa trên độ tương đồng vector. Hệ thống truy xuất các đoạn văn bản (text chunks) có bản nhúng vector gần nhất với câu hỏi mà không tận dụng cấu trúc đồ thị. Chế độ này phù hợp cho các câu hỏi đơn giản, yêu cầu trích xuất thông tin trực tiếp từ một đoạn văn cụ thể.

\indent Chế độ truy vấn cục bộ (Local Query): Chế độ này được tối ưu hóa cho các câu hỏi chi tiết về một khái niệm cụ thể (Low-level Query). Hệ thống sử dụng cơ chế tìm kiếm vector để xác định các thực thể liên quan nhất, sau đó mở rộng truy xuất các nút lân cận và các mối quan hệ trực tiếp trong đồ thị tri thức. Quy trình này cung cấp một ngữ cảnh tri thức tập trung, giúp mô hình ngôn ngữ trả lời chính xác các thuộc tính và mối quan hệ vi mô của thực thể.

\indent Chế độ truy vấn toàn cục (Global Query): Đối với các câu hỏi mang tính tổng quát cao (High-level Query), đòi hỏi sự hiểu biết về toàn bộ kho dữ liệu, chế độ toàn cục phát huy vai trò chủ đạo. Thay vì tìm kiếm đoạn văn rời rạc, hệ thống tận dụng cấu trúc đồ thị ở mức độ vĩ mô bằng cách tổng hợp các mô tả tóm tắt của thực thể và quan hệ quan trọng. Cách tiếp cận này giúp mô hình ngôn ngữ bao quát được các chủ đề lớn và đưa ra những nhận định mang tính hệ thống vượt qua giới hạn cửa sổ ngữ cảnh thông thường.

\indent Chế độ truy vấn hỗn hợp (Mix Query): Đây là sự kết hợp giữa tìm kiếm cục bộ và toàn cục. Chế độ Mix đồng thời truy xuất các chi tiết cụ thể từ đồ thị lân cận và các tóm tắt tri thức vĩ mô từ cấp độ toàn cục. Quy trình xử lý này đảm bảo câu trả lời vừa có độ chính xác của dữ liệu chi tiết, vừa có chiều sâu của lập luận hệ thống, đáp ứng tốt nhất cho các nhiệm vụ phân tích tri thức đa tầng.

Cơ chế truy xuất hai luồng cục bộ và toàn cục cùng quy trình lọc ngữ cảnh của LightRAG được thể hiện chi tiết ở Hình~\ref{fig:lightrag_retrieval_internal}.

 \begin{figure}[H]
	\centering
	\resizebox{\textwidth}{!}{
	\begin{tikzpicture}[
	    >=Stealth,
	    every node/.style={font=\sffamily\small},
	    storageBox/.style={rectangle, rounded corners=3pt, draw=black, fill=gray!15, thick, align=center, text width=4.8cm, minimum height=1.0cm},
	    nodeBox/.style={rectangle, rounded corners=3pt, draw=black!70, fill=white, thick, align=center, text width=3.8cm, minimum height=0.8cm},
	    inputBox/.style={rectangle, rounded corners=3pt, draw=black, fill=gray!30, thick, align=center, text width=3.2cm, minimum height=0.8cm},
	    outputBox/.style={rectangle, rounded corners=3pt, draw=black, fill=gray!30, thick, align=center, text width=3.2cm, minimum height=0.8cm},
	    labelTag/.style={font=\sffamily\scriptsize, fill=white, inner sep=1.5pt, rounded corners=2pt},
	    solidEdge/.style={-{Stealth[length=2.5mm]}, thick},
	]

	    % ==========================================
	    % 1. DINH VI CAC NODE
	    % ==========================================

	    %% STORAGE ROW (tren cung) -- gian rong +-14
	    \node[storageBox] (VDB_L) at (-14.0, 10.0) {Cơ sở dữ liệu Vector\\(Nano Vector DB - Trái)};
	    \node[storageBox] (KG)    at (  0.0, 10.0) {Đồ thị tri thức\\(Knowledge Graph)};
	    \node[storageBox] (VDB_R) at ( 14.0, 10.0) {Cơ sở dữ liệu Vector\\(Nano Vector DB - Phải)};

	    %% KHUNG RETRIEVE
	    \filldraw[fill=gray!3, draw=black!40, dashed, rounded corners=12pt, very thick] (-18.0, 6.8) rectangle (18.0, -1.1);
	    \node[font=\sffamily\bfseries\large, anchor=north west, color=black] at (-17.7, 6.5) {Retrieve (Phân hệ Truy xuất)};

	    %% TRUNG TAM: Keyword Extraction
	    \node[nodeBox] (KWEPrompt)    at (-5.8, -0.2) {Keywords Extraction Prompt\\(Mẫu gợi ý trích xuất)};
	    \node[nodeBox] (SysPromptExt) at (  0.0, -0.2) {System Prompt\\(Gợi ý hệ thống)};
	    \node[nodeBox] (KWExtract)    at (  0.0,  1.5) {keywords\_extraction\\(Trích xuất từ khóa)};
	    \node[nodeBox] (LowKW)        at ( -2.5,  3.0) {low\_level\_keywords\\(Từ khóa mức thấp)};
	    \node[nodeBox] (HighKW)       at (  2.5,  3.0) {high\_level\_keywords\\(Từ khóa mức cao)};
	    \node[nodeBox] (Embed1)       at ( -2.5,  4.8) {Embedding (Nhúng)};
	    \node[nodeBox] (Embed2)       at (  2.5,  4.8) {Embedding (Nhúng)};

	    %% NHANH TRAI (Low-Level Retrieval)
	    \node[nodeBox] (TopK_E)    at (-11.0,  5.5) {TopK Entities Results\\(Kết quả Top-K Thực thể)};
	    \node[nodeBox] (RelText_L) at (-16.0,  3.5) {Related text\_units\\(Đoạn văn bản liên quan)};
	    \node[nodeBox] (RelRel)    at ( -7.0,  3.5) {Related Relations\\(Quan hệ liên quan)};
	    \node[nodeBox] (LocalCtx)  at (-11.0,  1.5) {Local Query Context\\(Ngữ cảnh truy vấn cục bộ)};

	    %% NHANH PHAI (High-Level Retrieval)
	    \node[nodeBox] (TopK_R)    at ( 11.0,  5.5) {TopK Relations Results\\(Kết quả Top-K Quan hệ)};
	    \node[nodeBox] (RelText_R) at ( 16.0,  3.5) {Related text\_units\\(Đoạn văn bản liên quan)};
	    \node[nodeBox] (RelEnt)    at (  7.0,  3.5) {Related Entities\\(Thực thể liên quan)};
	    \node[nodeBox] (GlobalCtx) at ( 11.0,  1.5) {Global Query Context\\(Ngữ cảnh truy vấn toàn cục)};

	    %% QUERY INPUT
	    \node[inputBox] (Query) at (-17.0, -3.5) {Query\\(Câu hỏi)};

	    %% KHUNG GENERATION
	    \filldraw[fill=gray!3, draw=black!40, dashed, rounded corners=12pt, very thick] (-5.8, -1.4) rectangle (5.8, -5.5);
	    \node[font=\sffamily\bfseries\large, anchor=north west, color=black] at (-5.5, -1.7) {Generation (Tạo sinh)};
	    \node[nodeBox] (CombCtx)      at (  0.0, -2.6) {combined context\\(Ngữ cảnh kết hợp)};
	    \node[nodeBox] (SysTpl)       at ( -2.5, -4.4) {System Template Prompt\\(Mẫu gợi ý hệ thống)};
	    \node[nodeBox] (SysPromptGen) at (  2.5, -4.4) {System Prompt\\(Gợi ý hệ thống)};

	    %% RESPONSE OUTPUT
	    \node[outputBox] (Response) at (0.0, -7.2) {Response\\(Câu trả lời)};

	    % ==========================================
	    % 2. DINH TUYEN MUI TEN (goc vuong tuyet doi)
	    % Moi ket noi dung bus lane y rieng biet
	    % ==========================================

	    %% A. KWEPrompt -> SysPromptExt (ngang)
	    \draw[solidEdge] (KWEPrompt.east) -- (SysPromptExt.west);

	    %% B. SysPromptExt -> KWExtract (thang dung len)
	    \draw[solidEdge] (SysPromptExt.north) -- node[labelTag, text=black, fill=gray!15, pos=0.5] {use LLM} (KWExtract.south);

	    %% C. KWExtract -> LowKW (re trai tai y=2.2)
	    \draw[solidEdge] (KWExtract.north) -- (0.0, 2.2) -- (-2.5, 2.2) -- (LowKW.south);

	    %% D. KWExtract -> HighKW (re phai tu nut giao y=2.2)
	    \draw[solidEdge] (0.0, 2.2) -- (2.5, 2.2) -- (HighKW.south);

	    %% E. LowKW -> Embed1
	    \draw[solidEdge] (LowKW.north)  -- node[labelTag, text=black, fill=gray!15, pos=0.5] {use EB Model} (Embed1.south);

	    %% F. HighKW -> Embed2
	    \draw[solidEdge] (HighKW.north) -- node[labelTag, text=black, fill=gray!15, pos=0.5] {use EB Model} (Embed2.south);

	    %% G. Embed1 -> VDB_L (bus y=7.0 phia tren khung Retrieve)
	    \draw[solidEdge] (Embed1.north) -- (-2.5, 7.0) -- (-14.0, 7.0) -- (VDB_L.south);

	    %% H. Embed2 -> VDB_R (bus y=7.0, doi xung G)
	    \draw[solidEdge] (Embed2.north) -- ( 2.5, 7.0) -- ( 14.0, 7.0) -- (VDB_R.south);

	    %% I. VDB_L -> TopK_E  (bus y=9.0, lech +3mm, sang phai den x=-11.3)
	    \draw[solidEdge] ([xshift=3mm]VDB_L.south) -- (-13.7, 9.0) -- (-11.3, 9.0) -- ([xshift=-3mm]TopK_E.north);

	    %% J. VDB_L -> RelText_L  (bus y=8.6, lech -3mm, sang trai den x=-16.0)
	    \draw[solidEdge] ([xshift=-3mm]VDB_L.south) -- (-14.3, 8.6) -- (-16.0, 8.6) -- (RelText_L.north);

	    %% K. KG -> TopK_E  (bus y=9.3, lech -6mm, sang trai den x=-10.7)
	    \draw[solidEdge] ([xshift=-6mm]KG.south) -- (-0.6, 9.3) -- (-10.7, 9.3) -- ([xshift=3mm]TopK_E.north);

	    %% L. KG -> RelRel  (bus y=8.0, lech -2mm, sang trai den x=-7.0)
	    \draw[solidEdge] ([xshift=-2mm]KG.south) -- (-0.2, 8.0) -- (-7.0, 8.0) -- (RelRel.north);

	    %% M. VDB_R -> TopK_R  (bus y=9.0, doi xung I)
	    \draw[solidEdge] ([xshift=-3mm]VDB_R.south) -- (13.7, 9.0) -- (11.3, 9.0) -- ([xshift=3mm]TopK_R.north);

	    %% N. VDB_R -> RelText_R  (bus y=8.6, doi xung J)
	    \draw[solidEdge] ([xshift=3mm]VDB_R.south) -- (14.3, 8.6) -- (16.0, 8.6) -- (RelText_R.north);

	    %% O. KG -> TopK_R  (bus y=9.3, doi xung K)
	    \draw[solidEdge] ([xshift=6mm]KG.south) -- (0.6, 9.3) -- (10.7, 9.3) -- ([xshift=-3mm]TopK_R.north);

	    %% P. KG -> RelEnt  (bus y=8.0, doi xung L)
	    \draw[solidEdge] ([xshift=2mm]KG.south) -- (0.2, 8.0) -- (7.0, 8.0) -- (RelEnt.north);

	    %% Q. TopK_E -> LocalCtx  (thang dung xuong)
	    \draw[solidEdge] (TopK_E.south) -- (LocalCtx.north);

	    %% R. RelText_L -> LocalCtx  (bus y=2.8, lech -5mm o LocalCtx)
	    \draw[solidEdge] (RelText_L.south) -- (-16.0, 2.8) -- (-11.5, 2.8) -- ([xshift=-5mm]LocalCtx.north);

	    %% S. RelRel -> LocalCtx  (bus y=2.6, lech +5mm o LocalCtx)
	    \draw[solidEdge] (RelRel.south) -- (-7.0, 2.6) -- (-10.5, 2.6) -- ([xshift=5mm]LocalCtx.north);

	    %% T. TopK_R -> GlobalCtx  (thang dung xuong, doi xung Q)
	    \draw[solidEdge] (TopK_R.south) -- (GlobalCtx.north);

	    %% U. RelText_R -> GlobalCtx  (bus y=2.8, doi xung R)
	    \draw[solidEdge] (RelText_R.south) -- (16.0, 2.8) -- (11.5, 2.8) -- ([xshift=5mm]GlobalCtx.north);

	    %% V. RelEnt -> GlobalCtx  (bus y=2.6, doi xung S)
	    \draw[solidEdge] (RelEnt.south) -- (7.0, 2.6) -- (10.5, 2.6) -- ([xshift=-5mm]GlobalCtx.north);

	    %% W. LocalCtx -> CombCtx  (xuong den y=-2.6, roi sang phai)
	    \draw[solidEdge] (LocalCtx.south) -- (-11.0, -2.6) -- (CombCtx.west);

	    %% X. GlobalCtx -> CombCtx  (xuong den y=-2.6, roi sang trai)
	    \draw[solidEdge] (GlobalCtx.south) -- (11.0, -2.6) -- (CombCtx.east);

	    %% Y. CombCtx -> SysPromptGen  (xuong den y=-3.6, roi sang phai)
	    \draw[solidEdge] (CombCtx.south) -- (0.0, -3.6) -- (2.5, -3.6) -- (SysPromptGen.north);

	    %% Z. SysTpl -> SysPromptGen  (ngang)
	    \draw[solidEdge] (SysTpl.east) -- (SysPromptGen.west);

	    %% AA. Query -> SysPromptExt  (luon qua khe y=-1.25 giua 2 khung)
	    \draw[solidEdge] (Query.north) -- (-17.0, -1.25) -- (0.0, -1.25) -- (SysPromptExt.south);

	    %% AB. Query -> SysPromptGen  (luon qua day y=-6.3)
	    \draw[solidEdge] (Query.south) -- (-17.0, -6.3) -- (2.5, -6.3) -- (SysPromptGen.south);

	    %% AC. SysPromptGen -> Response  (ra phai x=5.3, xuong y=-6.7, vao Response)
	    \draw[solidEdge] (SysPromptGen.east) -- (5.3, -4.4) -- (5.3, -6.7) -- (0.0, -6.7) -- (Response.north);

	\end{tikzpicture}
	}
	\caption{Cơ chế truy xuất tri thức và lọc ngữ cảnh hai tầng trong LightRAG}
	\label{fig:lightrag_retrieval_internal}
 \end{figure}

Việc viết lại câu hỏi thành dạng độc lập kết hợp với từ khóa giúp khắc phục hiện tượng mất ngữ cảnh khi câu hỏi của sinh viên quá ngắn hoặc mơ hồ. Toàn bộ trạng thái hội thoại được lưu bền vững qua cơ chế checkpoint trên MongoDB, cho phép khôi phục đúng bước đang xử lý khi có gián đoạn. Để gia tăng độ linh hoạt trong việc đáp ứng các nhu cầu tra cứu đa dạng, hệ thống cung cấp hai cấu hình truy vấn chuyên biệt bao gồm cấu hình truy vấn "flash" và cấu hình truy vấn "pro". Cấu hình "flash" ưu tiên tối ưu hóa thời gian phản hồi bằng cách giới hạn số lượng kết quả trả về từ cơ sở dữ liệu, trong khi cấu hình "pro" tập trung khai thác triệt để các mối quan hệ phức tạp nhằm mang lại ngữ cảnh toàn diện nhất. Song song với đó, hệ thống tích hợp thêm một luồng tra cứu đồ thị trực tiếp bằng ngôn ngữ truy vấn cấu trúc. Luồng xử lý này cho phép hệ thống đếm số lượng thực thể hoặc kiểm chứng các sự kiện cụ thể một cách độc lập mà không cần phải ánh xạ qua không gian vector nhúng, qua đó giảm thiểu đáng kể chi phí tính toán. Đối với các truy vấn đòi hỏi thông tin ngữ nghĩa sâu hơn thông qua các chế độ tìm kiếm hỗn hợp, sự chính xác của ngữ cảnh được đảm bảo bởi một cơ chế lọc kết quả nghiêm ngặt ngay sau khi quá trình truy xuất từ cơ sở dữ liệu hoàn tất. Dựa trên các thông số đặc tả như ngành học hay học kỳ đã được trích xuất từ bước phân tích ý định, hệ thống tự động đối chiếu và loại bỏ các mảnh thông tin nhiễu hoặc sai lệch bối cảnh. Quá trình này đảm bảo kho liệu đầu vào cho bước sinh văn bản hoàn toàn tinh sạch. Việc chuẩn bị ngữ cảnh còn được cá nhân hóa sâu sắc thông qua việc lồng ghép trực tiếp hồ sơ bộ nhớ dài hạn của người dùng vào trong hướng dẫn tạo sinh, giúp câu trả lời gắn liền với định hướng của từng sinh viên. Nhằm đảm bảo tính chuẩn xác tuyệt đối trước khi phản hồi người dùng, kiến trúc điều phối tích hợp một tác tử kiểm chứng học thuật hoạt động như một màng lọc cuối cùng. Tác tử này được thiết kế chuyên biệt để thẩm định tính xác thực đối với các nhóm câu hỏi hệ trọng như lộ trình học tập, môn tiên quyết hay quy định tín chỉ. Quá trình kiểm chứng liên tục đối chiếu chéo câu trả lời dự kiến với các quy chuẩn học thuật cốt lõi nhằm ngăn chặn tình trạng ảo giác thông tin của mô hình tạo sinh. Cuối cùng, hệ thống áp dụng cơ chế xử lý bất đồng bộ toàn diện, cho phép truyền phát dữ liệu theo dạng chuỗi luồng liên tục ngay khi có kết quả, đồng thời thực thi các tác vụ cập nhật bộ nhớ dài hạn dưới nền sau khi kết thúc phản hồi, qua đó mang lại trải nghiệm tương tác liền mạch và tối ưu nhất.

\subsection{Thuật toán điều phối (Pseudocode)}

\begin{algorithm}[H]
	\caption{Quy trình điều phối Agent}\label{alg:orchestration}
	\begin{algorithmic}[1]
		\Require $query,\ user\_id$ \Comment{Đầu vào từ sinh viên}
		\Ensure $answer$ \Comment{Câu trả lời cuối cùng}

		\State $routing \gets \text{AnalyzeAgent}(query,\ \text{Summarize}(history))$
		\State $decision \gets routing.decision$ \Comment{direct\_answer $\mid$ retrieve $\mid$ reject}
		\If{$decision = \text{``reject''}$}
			\State \Return $\text{REJECT\_MESSAGE}$
		\EndIf
		\State $memory \gets \text{LoadFromMongo}(user\_id)$ \Comment{chỉ khi $routing.needs\_lt$}
		\State $genReq \gets \text{BuildGenerationRequest}(routing,\ memory)$
		\If{$decision = \text{``retrieve''}$}
			\State $answer \gets \text{LightRAG}(genReq.query,\ genReq.mode,\ K_{high},\ K_{low})$
		\Else
			\State $answer \gets \text{LLM}(genReq.system\_prompt,\ query)$
		\EndIf
		\If{$routing.should\_update\_lt$}
			\State $\text{UpdateMemory}(user\_id,\ query,\ answer)$ \Comment{trích xuất + hợp nhất hồ sơ}
		\EndIf
		\State \Return $answer$
	\end{algorithmic}
\end{algorithm}

\section{Thiết lập thực nghiệm}

\subsection{Thiết lập thực nghiệm phân hệ trích xuất (MinerU)}

\subsubsection{Môi trường thực nghiệm}
Hệ thống được triển khai thực nghiệm trên hạ tầng máy chủ đám mây nhằm đảm bảo khả năng tính toán mạnh mẽ cho các tác vụ xử lý ảnh và văn bản. Quá trình xử lý tài liệu đầu vào sử dụng công cụ MinerU (từ phiên bản 3.2.3) tích hợp các mô hình Vision-Language Model (VLM), OCR và Mathematical Formula Recognition (MFR). 

\subsubsection{Mô tả dữ liệu và quá trình thu thập}
Để thực nghiệm và đánh giá chất lượng của bộ giải mã MinerU, nghiên cứu đã tiến hành thu thập tập dữ liệu Chương trình đào tạo chuẩn của sáu ngành và khối ngành thuộc Khoa, bao gồm: Cử nhân tài năng, Công nghệ thông tin, Hệ thống thông tin, Khoa học máy tính, Kỹ thuật phần mềm, và Trí tuệ nhân tạo. Nhằm đảm bảo tính đại diện, tài liệu được phân rã theo từng khóa tuyển sinh từ năm 2022 đến năm 2025. 

\subsubsection{Độ đo để đánh giá (Performance Metrics)}
Phương pháp đánh giá được thực hiện thông qua việc đối sánh kết quả trích xuất tự động với nhãn chuẩn (Ground Truth). Để đánh giá mức độ chính xác cơ bản, nghiên cứu sử dụng CER (Character Error Rate) và WER (Word Error Rate) nhằm phản ánh tỷ lệ sai sót ở cấp độ ký tự và từ. Tương tự, Edit Distance Sim (Độ tương đồng khoảng cách sửa đổi) được áp dụng để đo lường sự sai khác chuỗi tổng thể. Nhằm định lượng mức độ trùng khớp văn bản sâu hơn, BLEU và METEOR được sử dụng để đánh giá độ bảo toàn n-gram. Đặc biệt, đối với cấu trúc dữ liệu phức tạp, TEDS-Struct và TEDS-Content (Tree-Edit-Distance-based Similarity) đóng vai trò đo lường chuyên biệt cho mức độ bảo toàn cấu trúc bảng biểu. 

Biện luận về độ quan trọng của các chỉ số: Đối với kiến trúc truy xuất tăng cường đồ thị, các chỉ số BLEU, TEDS-Struct, TEDS-Content và Edit Distance Sim đóng vai trò then chốt nhất. Nếu mô hình bóc tách thất bại trong việc tái tạo cấu trúc bảng (TEDS-Struct thấp), lượng tri thức sẽ bị suy biến. Chỉ số BLEU cao minh chứng văn bản bảo toàn chính xác các cụm từ gốc, giúp LLM không sai lệch ngữ cảnh. Edit Distance Sim đảm bảo tính nguyên vẹn của mã học phần, tạo tiền đề để thuật toán đối soát và hợp nhất thực thể hoạt động ổn định.

\subsection{Thiết lập thực nghiệm Hệ thống Đồ thị tri thức và Đa tác nhân}
\subsubsection{Môi trường thực nghiệm}
Đối với quá trình đánh giá chất lượng đồ thị tri thức bằng phương pháp LLM-as-a-Judge, nhóm nghiên cứu lựa chọn mô hình Qwen2.5-72B-Instruct \cite{qwen2.5} làm trọng tài đánh giá (Judge). Việc lựa chọn mô hình này dựa trên ba lý do chính. Trước tiên, Qwen2.5-72B-Instruct là dòng mô hình mã nguồn mở hàng đầu với 72 tỷ tham số, sở hữu năng lực tư duy logic và lập luận phức tạp vượt trội so với các mô hình nhỏ hơn, phù hợp với yêu cầu đánh giá ngữ nghĩa sâu của các quan hệ trong đồ thị. Tiếp theo, mô hình này có khả năng hiểu và phân tích tiếng Việt chuyên ngành xuất sắc - điều này đã được xác thực khoa học qua bộ đánh giá VMLU Benchmarks (2025), trong đó Qwen2.5-72B-Instruct đạt hiệu năng hàng đầu trong hầu hết các tác vụ ngôn ngữ tiếng Việt \cite{vmlu2025}, giúp phán xét chính xác các khái niệm đặc thù của giáo dục đại học Việt Nam.

\subsubsection{Mô tả dữ liệu và quá trình thu thập}
Tiếp nối quá trình tiền xử lý văn bản bằng công cụ MinerU, tập dữ liệu đầu vào của hệ thống được xây dựng dựa trên các tài liệu thực tế về chương trình đào tạo, đề cương môn học và các quy chế học vụ của Khoa Công nghệ Thông tin, Trường Đại học Khoa học Tự nhiên, ĐHQG-HCM. Tập dữ liệu bao gồm nhiều tài liệu chính, được chia thành ba nhóm chính. Nhóm thứ nhất là các mô tả chương trình đào tạo chuẩn của sáu ngành và khối ngành thuộc Khoa, gồm: Cử nhân tài năng, Công nghệ thông tin, Hệ thống thông tin, Khoa học máy tính, Kỹ thuật phần mềm và Trí tuệ nhân tạo, được phân rã theo từng khóa tuyển sinh từ năm 2022 đến năm 2025. Nhóm thứ hai tập trung vào đề cương chi tiết của hàng loạt học phần từ đại cương đến chuyên ngành, chứa đựng các thông tin cốt lõi như chuẩn đầu ra (ELO) và điều kiện tiên quyết. Nhóm thứ ba bao gồm các văn bản quy định và thông báo học vụ, điển hình là \textit{Quy định chuẩn đầu ra Tiếng Anh từ Khóa 2022} và \textit{Thông báo đăng ký học phần Học kỳ 2 Năm học 2025-2026}. Việc lựa chọn tập dữ liệu nội bộ và có tính đặc thù cao nhằm mục đích kiểm chứng khả năng của hệ thống LightRAG trong việc thấu hiểu và bóc tách các tri thức chuyên ngành phức tạp, vốn đòi hỏi sự liên kết chặt chẽ về mặt logic.
\subsubsection{Đánh giá chất lượng nội tại của đồ thị tri thức (Graph Quality Metrics)}

Nghiên cứu này chú trọng vào việc đo lường chất lượng nội tại của đồ thị (intrinsic graph quality) dựa trên mức độ trung thực so với văn bản gốc. Để thực hiện điều này, nghiên cứu áp dụng tổ hợp ba phương pháp đánh giá tiên tiến nhằm giải quyết ba vấn đề cốt lõi thường gặp khi trích xuất đồ thị tự động: tình trạng phân mảnh, hiện tượng ảo giác, và lỗi thiết kế lược đồ. 

Đầu tiên là, nhằm giải quyết tình trạng phân mảnh thông tin - khi đồ thị tạo ra các thực thể rời rạc thay vì một mạng lưới liên kết, nghiên cứu áp dụng phương pháp đánh giá tính toàn vẹn cấu trúc (Graph Structure Health) theo đề xuất của Seo và cộng sự \cite{seo2022structural}. Cụ thể, phương pháp này đo lường số lượng các nút mồ côi (dangling nodes) và thống kê tỷ lệ bao phủ văn bản (chunk coverage). Nhờ đó, hệ thống có thể đảm bảo rằng hầu hết các thực thể trích xuất ra đều được kết nối chặt chẽ và không bỏ sót thông tin quan trọng từ tài liệu gốc. Kế tiếp để giảm thiểu hiện tượng ảo giác (hallucination) - tức tình trạng mô hình ngôn ngữ tự bịa đặt ra các mối quan hệ sai lệch (ví dụ: tự sáng tác ra môn học tiên quyết không tồn tại) - nghiên cứu sử dụng phương pháp đánh giá LLM-as-a-Judge do Edge và cộng sự (nhóm tác giả GraphRAG) \cite{edge2024} khởi xướng. Áp dụng phương pháp này, nghiên cứu lấy mẫu ngẫu nhiên các quan hệ sinh ra và yêu cầu mô hình Qwen2.5-72B đóng vai trò giám khảo độc lập. Mô hình sẽ đối chiếu trực tiếp từng mối quan hệ với đoạn văn bản gốc tương ứng để chấm điểm tính có căn cứ (groundedness). Phương pháp này giúp hệ thống cam kết độ tin cậy tuyệt đối của tri thức được trích xuất. Sau cùng là ngăn chặn tình trạng trôi dạt lược đồ (schema drift) và đảm bảo cấu trúc đồ thị thực sự phục vụ được mục tiêu truy vấn, nghiên cứu áp dụng chuẩn kiểm định tính hợp lệ lược đồ (Ontology Validation). Theo hướng tiếp cận của hệ thống OOPS! do Poveda-Villalón và cộng sự \cite{poveda2014oops} phát triển, kết hợp với các tiêu chuẩn thiết kế siêu dữ liệu của LightRAG \cite{guo2025lightrag}, phương pháp này tự động rà soát các lỗi logic trong thiết kế đồ thị và kiểm tra khả năng trả lời các câu hỏi năng lực (Competency Questions). Việc này giúp phát hiện sớm các cấu trúc đồ thị lỗi từ trong trứng nước trước khi đưa vào truy vấn thực tế. Về mặt phương pháp luận, quy trình kiểm định tính hợp lệ dựa trên danh mục bẫy lỗi (pitfalls) của OOPS! vẫn là một chuẩn mực cốt lõi được áp dụng xuyên suốt trong các nghiên cứu hiện đại. Minh chứng là trong một nghiên cứu y tế công bố năm 2026 \cite{franco2026meso}, phương pháp validation này đã được sử dụng thành công để kiểm thử tự động một hệ thống phân cấp tri thức y học phức tạp (MeSO). Bằng cách quét và phát hiện các lỗi cấu trúc như thiếu ràng buộc thuộc tính nghịch đảo hoặc sai lệch phân cấp, thực nghiệm năm 2026 đã chứng minh tính hiệu quả vượt thời gian của phương pháp này trong việc duy trì sự nhất quán logic và bảo vệ tính toàn vẹn của các mô hình đồ thị tri thức trước khi đưa vào vận hành.

\subsubsection{Đánh giá chất lượng sinh văn bản và hiệu năng điều phối (RAG \& Agentic Metrics)}

Bên cạnh chất lượng đồ thị, chất lượng truy xuất và sinh văn bản cũng được đánh giá nghiêm ngặt thông qua khung đo lường RAGAS (Retrieval-Augmented Generation Assessment) \cite{es2023ragas}. Nghiên cứu áp dụng ba chỉ số cốt lõi của phương pháp này. Đầu tiên, độ liên quan của câu trả lời (Answer Relevance) được tính toán nhằm đảm bảo phản hồi đi thẳng vào trọng tâm truy vấn. Thứ hai, độ trung thực (Faithfulness) đo lường khả năng sinh văn bản bám sát ngữ cảnh từ đồ thị, đóng vai trò là lớp phòng thủ thứ hai giúp hạn chế tối đa hiện tượng ảo giác. Cuối cùng, độ gọi bối cảnh (Context Recall) đánh giá năng lực của hệ thống trong việc thu thập đủ các khái niệm phân tán trên đồ thị để tạo nên một câu trả lời hoàn chỉnh.

Về mặt hiệu năng kiến trúc, nghiên cứu tiến hành đo lường năng lực điều phối đa tác nhân (Multi-Agent). Cụ thể, tỷ lệ điều phối đúng (Routing Accuracy) được sử dụng để đánh giá độ chính xác của Tác tử phân tích (Orchestrator) khi quyết định phân luồng câu hỏi vào các chế độ xử lý phù hợp (Local, Global, Hybrid hay Cypher Lookup). Đồng thời, tốc độ phản hồi trung bình (Average Latency) được ghi nhận nhằm so sánh độ trễ giữa luồng xử lý đồ thị trực tiếp so với luồng RAG truyền thống, từ đó chứng minh hiệu quả tối ưu hóa thời gian thực của kiến trúc phân tuyến mới.









\clearpage
